\newMonth{Junius}
\setrunningtitles{Junius}{June}

% ---- martyrology/mart06/mart0601.htm
\needspace{10\baselineskip}
\begin{paracol}{2}
\selectlanguage{latin}
\begin{center}{\color{gregoriocolor} Kaléndis Júnii. 
 Luna\dots\ }\end{center}
\switchcolumn
\selectlanguage{english}
\begin{center}{\color{gregoriocolor} The 1\textsuperscript{st} Day of 
 June. The\dots\ Day of the Moon.}\end{center}
\end{paracol}

\noindent\begin{tabularx}{\linewidth}{*{19}{>{\centering\arraybackslash}X}}
 \textcolor{gregoriocolor}{a} & \textcolor{gregoriocolor}{b} & \textcolor{gregoriocolor}{c} & \textcolor{gregoriocolor}{d} & \textcolor{gregoriocolor}{e} & \textcolor{gregoriocolor}{f} & \textcolor{gregoriocolor}{g} & \textcolor{gregoriocolor}{h} & \textcolor{gregoriocolor}{i} & \textcolor{gregoriocolor}{k} & \textcolor{gregoriocolor}{l} & \textcolor{gregoriocolor}{m} & \textcolor{gregoriocolor}{n} & \textcolor{gregoriocolor}{p} & \textcolor{gregoriocolor}{q} & \textcolor{gregoriocolor}{r} & \textcolor{gregoriocolor}{s} & \textcolor{gregoriocolor}{t} & \textcolor{gregoriocolor}{u} \\
 5 & 6 & 7 & 8 & 9 & 10 & 11 & 12 & 13 & 14 & 15 & 16 & 17 & 18 & 19 & 20 & 21 & 22 & 23 \\
\end{tabularx}
\vspace{0.5\baselineskip}
\noindent\begin{tabularx}{\linewidth}{*{12}{>{\centering\arraybackslash}X}}
 \textcolor{gregoriocolor}{A} & \textcolor{gregoriocolor}{B} & \textcolor{gregoriocolor}{C} & \textcolor{gregoriocolor}{D} & \textcolor{gregoriocolor}{E} & F & \textcolor{gregoriocolor}{F} & \textcolor{gregoriocolor}{G} & \textcolor{gregoriocolor}{H} & \textcolor{gregoriocolor}{M} & \textcolor{gregoriocolor}{N} & \textcolor{gregoriocolor}{P} \\
 24 & 25 & 26 & 27 & 28 & 29 & 29 & 30 & 1 & 2 & 3 & 4 \\
\end{tabularx}

\begin{paracol}{2}
\selectlanguage{latin}
\lettrine[lines=2]{S}{anctæ} Angelæ Meríci 
 Vírginis, ex tértio Ordine sancti Francísci; quæ fuit Institútrix Societátis 
 Vírginum sanctæ Ursulæ, et ad accipiéndam corónam immarcescíbilem, sexto 
 Kaléndas Februárii a cælésti Sponso vocáta est.
\switchcolumn
\selectlanguage{english}
\lettrine[lines=2]{S}{t.} Angela Merici, virgin of the 
 Third Order of St. Francis. She was the foundress of the Nuns of St. 
 Ursula, and was called by her heavenly Spouse on the 27th of January in 
 order to receive an incorruptible crown.
\switchcolumn*
\selectlanguage{latin}
Romæ sancti Juvéntii 
 Mártyris.
\switchcolumn
\selectlanguage{english}
At Rome, St. Juventius, martyr.
\switchcolumn*
\selectlanguage{latin}
Augustodúni sanctórum 
 Reveriáni Epíscopi, et Pauli Presbyteri, cum áliis decem, qui, sub Aureliáno 
 Príncipe, martyrio coronáti sunt.
\switchcolumn
\selectlanguage{english}
At Autun, the Saints Reverian, 
 bishop, and Paul, a priest, along with ten others, who were crowned with 
 martyrdom under Emperor Aurelian.
\switchcolumn*
\selectlanguage{latin}
Cæsaréæ, in Palæstína, 
 beáti Pámphili, Presbyteri et Mártyris, viri admirándæ sanctitátis et 
 doctrínæ, atque in páuperes munífici; qui, ob Christi fidem, in persecutióne 
 Galérii Maximiáni, primum, sub Urbáno Præside, cruciátus et in cárcerem 
 trusus, deínde, sub Firmiliáno, íterum revocátus ad pœnas, una cum áliis 
 martyrium consummávit. Passi sunt étiam tunc Valens Diáconus, et 
 Paulus, aliíque novem; quorum memória áliis diébus celebrátur.
\switchcolumn
\selectlanguage{english}
At Caesarea in Palestine, blessed 
 Pamphilus, priest and martyr, a man of remarkable sanctity and learning, and 
 great charity to the poor. In the persecution of Galerius Maximian, he 
 was tortured for the faith of Christ, under Governor Urbanus, and thrown 
 into prison. Later he was again subjected to torments under Firmilian, 
 and he completed his martyrdom with others. At the same time, there 
 suffered Valens, a deacon, and Paul, and nine others, whose commemoration 
 occurs on other days.
\switchcolumn*
\selectlanguage{latin}
In Cappadócia sancti 
 Thespésii Mártyris, qui, sub Alexándro Imperatóre et Simplício Præfécto, post 
 ália torménta, decollátus est.
\switchcolumn
\selectlanguage{english}
In Cappadocia, in the time of 
 Emperor Alexander and the prefect Simplicius, the holy martyr Thespesius, 
 who, after undergoing many torments, was beheaded.
\switchcolumn*
\selectlanguage{latin}
In Ægypto sanctórum 
 Mártyrum Ischyriónis, ductóris mílitum, et aliórum quinque mílitum; qui, sub 
 Diocletiáno Imperatóre, pro fide Christi, divérso mortis génere perémpti 
 sunt.
\switchcolumn
\selectlanguage{english}
In Egypt, under Emperor Diocletian, 
 the holy martyrs Ischyrion, a military officer, and five other soldiers, who 
 were put to death in various ways for the faith of Christ.
\switchcolumn*
\selectlanguage{latin}
Item sancti Firmi 
 Mártyris, qui, in persecutióne Maximiáni, acerbíssimis plagis afféctus, 
 lapídibus percússus, ac demum cápite cæsus est.
\switchcolumn
\selectlanguage{english}
Also, St. Firmus, martyr, who was 
 scourged most severely, struck with stones, and finally beheaded during the 
 persecution of Maximian.
\switchcolumn*
\selectlanguage{latin}
Perúsiæ sanctórum 
 Mártyrum Felíni et Gratiniáni mílitum, qui, sub Décio, váriis torméntis 
 cruciáti, martyrii palmam gloriósa morte percepérunt.
\switchcolumn
\selectlanguage{english}
At Perugia, the holy martyrs 
 Felinus and Gratinian, soldiers under Decius, who were tortured in several 
 ways, and by a glorious death won the palm of martyrdom.
\switchcolumn*
\selectlanguage{latin}
Bonóniæ sancti Próculi 
 Mártyris, qui sub Maximiáno Imperatóre passus est.
\switchcolumn
\selectlanguage{english}
At Bologna, St. Proculus, martyr, 
 who suffered under Emperor Maximian.
\switchcolumn*
\selectlanguage{latin}
Amériæ, in Umbria, 
 sancti Secúndi Mártyris, qui, sub Diocletiáno, in Tíberim projéctus, 
 martyrium consummávit.
\switchcolumn
\selectlanguage{english}
At Amelia in Umbria, in the reign 
 of Diocletian, St. Secundus, martyr, who fulfilled his martyrdom when thrown 
 into the Tiber.
\switchcolumn*
\selectlanguage{latin}
Apud Tiférnum, in 
 Umbria, sancti Crescentiáni, mílitis Románi, qui sub eódem Imperatóre, 
 martyrio coronátus est.
\switchcolumn
\selectlanguage{english}
At Tiferno in Umbria, St. 
 Crescentian, a Roman soldier, crowned with martyrdom under the same emperor.
\switchcolumn*
\selectlanguage{latin}
In monastério 
 Lirinénsi, in Gállia, sancti Caprásii Abbátis.
\switchcolumn
\selectlanguage{english}
In the monastery of Lerins, the 
 abbot St. Caprasius.
\switchcolumn*
\selectlanguage{latin}
In monastério Onniénsi, 
 apud Burgos, in Hispánia, sancti Enecónis, Abbátis Benedictíni, ob 
 sanctitátis et miraculórum glóriam illústris.
\switchcolumn
\selectlanguage{english}
At Burgos in Spain, in the 
 monastery of Onia, St. Eneco, Benedictine abbot, made illustrious by his 
 sanctity and miracles.
\switchcolumn*
\selectlanguage{latin}
Apud Montem Falcum, in 
 Umbria, sancti Fortunáti Presbyteri, virtútibus et miráculis clari.
\switchcolumn
\selectlanguage{english}
At Montefalco in Umbria, St. 
 Fortunatus, a priest renowned for his virtues and his miracles.
\switchcolumn*
\selectlanguage{latin}
Tréviris sancti 
 Simeónis Mónachi, qui a Benedícto Papa Nono in Sanctórum númerum relátus 
 est.
\switchcolumn
\selectlanguage{english}
At Treves, St. Simeon, a monk, whom 
 Pope Benedict IX numbered among the saints.
\switchcolumn*
\selectlanguage{latin}
\end{paracol}

% \continueMonth{Junius}


% ---- martyrology/mart06/mart0602.htm
\needspace{10\baselineskip}
\begin{paracol}{2}
\selectlanguage{latin}
\begin{center}{\color{gregoriocolor} Quarto Nonas Júnii. 
 Luna\dots\ }\end{center}
\switchcolumn
\selectlanguage{english}
\begin{center}{\color{gregoriocolor} The 2\textsuperscript{nd} Day of 
 June. The\dots\ Day of the Moon.}\end{center}
\end{paracol}

\noindent\begin{tabularx}{\linewidth}{*{19}{>{\centering\arraybackslash}X}}
 \textcolor{gregoriocolor}{a} & \textcolor{gregoriocolor}{b} & \textcolor{gregoriocolor}{c} & \textcolor{gregoriocolor}{d} & \textcolor{gregoriocolor}{e} & \textcolor{gregoriocolor}{f} & \textcolor{gregoriocolor}{g} & \textcolor{gregoriocolor}{h} & \textcolor{gregoriocolor}{i} & \textcolor{gregoriocolor}{k} & \textcolor{gregoriocolor}{l} & \textcolor{gregoriocolor}{m} & \textcolor{gregoriocolor}{n} & \textcolor{gregoriocolor}{p} & \textcolor{gregoriocolor}{q} & \textcolor{gregoriocolor}{r} & \textcolor{gregoriocolor}{s} & \textcolor{gregoriocolor}{t} & \textcolor{gregoriocolor}{u} \\
 6 & 7 & 8 & 9 & 10 & 11 & 12 & 13 & 14 & 15 & 16 & 17 & 18 & 19 & 20 & 21 & 22 & 23 & 24 \\
\end{tabularx}
\vspace{0.5\baselineskip}
\noindent\begin{tabularx}{\linewidth}{*{12}{>{\centering\arraybackslash}X}}
 \textcolor{gregoriocolor}{A} & \textcolor{gregoriocolor}{B} & \textcolor{gregoriocolor}{C} & \textcolor{gregoriocolor}{D} & \textcolor{gregoriocolor}{E} & F & \textcolor{gregoriocolor}{F} & \textcolor{gregoriocolor}{G} & \textcolor{gregoriocolor}{H} & \textcolor{gregoriocolor}{M} & \textcolor{gregoriocolor}{N} & \textcolor{gregoriocolor}{P} \\
 25 & 26 & 27 & 28 & 29 & 1 & 30 & 1 & 2 & 3 & 4 & 5 \\
\end{tabularx}

\begin{paracol}{2}
\selectlanguage{latin}
\lettrine[lines=2]{R}{omæ} natális sanctórum 
 Mártyrum Marcellíni Presbyteri, et Petri Exorcístæ, qui, sub Diocletiáno, 
 cum multos in cárcere ad fidem erudírent, ídeo post dira víncula et plúrima 
 torménta, a Seréno Júdice decolláti sunt in loco qui dicebátur Silva Nigra; 
 quæ deínde, in honórem Sanctórum, commutáto nómine Silva Cándida appelláta 
 est. Horum córpora in crypta, juxta sanctum Tibúrtium, sepúlta sunt; 
 eorúmque sepúlcrum sanctus Dámasus Papa vérsibus póstea exornávit.
\switchcolumn
\selectlanguage{english}
\lettrine[lines=2]{A}{t} Rome, the birthday of the holy 
 martyr Marcellinus, priest, and Peter, exorcist, who instructed in the faith 
 many persons kept in prison. Under Diocletian, they were loaded with 
 chains, and after enduring many torments, were beheaded by Judge Serenus, in 
 a place which was then called the Black Forest, but which was in their 
 honour afterwards known as the White Forest. Their bodies were buried 
 in a crypt near St. Tiburtius, and Pope St. Damasus composed an epitaph in 
 verse for their tomb.
\switchcolumn*
\selectlanguage{latin}
In Campánia sancti 
 Erásmi, Epíscopi et Mártyris. Hic, sub Diocletiáno Augústo, primum 
 plumbátis cæsus, deínde fústibus gráviter mactátus, post resína, súlphure, 
 plumbo, pice, cera oleóque perfúsus, illæsus appáruit; deínde Fórmiis, sub 
 Maximiáno, divérsis atque immaníssimis supplíciis íterum tortus, sed ad 
 confirmándum céteros a Deo servátus; tandem, vocánte Dómino, martyrio clarus, 
 sancto fine quiévit. Ipsíus autem corpus Cajétam póstea translátum 
 est.
\switchcolumn
\selectlanguage{english}
In Campania, during the reign of 
 Decius, St. Erasmus, bishop and martyr, who was first scourged with leaded 
 whips and then severely beaten with rods. He also had resin, 
 brimstone, lead, pitch, wax, and oil poured over him, without receiving any 
 injury. Afterwards, under Maximian, he was again subjected to various 
 and most horrible tortures at Mola, but still was preserved from death by 
 the power of God in order to confirm others in the faith. Finally, 
 celebrated for his sufferings, and called by God, he closed his life by a 
 peaceful and holy death. His body was afterwards transferred to Gaeta.
\switchcolumn*
\selectlanguage{latin}
Lugdúni, in Gállia, 
 sanctórum Mártyrum Pothíni Epíscopi, Sancti Diáconi, Vétii, Epágathi, Matúri, 
 Póntici, Bíblidis, Attali, Alexándri et Blandínæ, cum áliis multis; quorum 
 fórtia et iteráta certámina, témpore Marci Aurélii Antoníni et Lúcii Veri, 
 Ecclésiæ Lugdunénsis epístola, ad Ecclésias Asiæ et Phrygiæ scripta, 
 recénset. In his sancta Blandína, sexu infírmior, córpore imbecíllior, 
 conditióne dejéctior, diuturnióra et acerbióra certámina súbiens, et fortis 
 adhuc pérmanens, gládio juguláta, céteros secúta est, quos hortabátur ad 
 palmam.
\switchcolumn
\selectlanguage{english}
At Lyons, many holy martyrs (Photinus, 
 a bishop, Sanctus, a deacon, Vetius, Epagathus, Maturus, Ponticus, Biblis, 
 Attalus, Alexander, and Blandina, with many others), whose many valiant 
 trials in the time of Marcus Aurelius Antoninus and Lucius Verus are 
 recorded in a letter from the church at Lyons to the churches of Asia and 
 Phrygia. St. Blandina, one of these martyrs, was weaker by reason of 
 her sex, more infirm in body, and of a lower station in life, and yet she 
 encountered longer and more terrible trials than the rest. But 
 remaining unshaken, she was put to the sword, and followed those whom she 
 had exhorted to win the palm of martyrdom.
\switchcolumn*
\selectlanguage{latin}
In ínsula Proconnéso, 
 in Propóntide, sancti Nicéphori, Epíscopi Constantinopolitáni, qui, 
 paternárum traditiónum propugnátor acérrimus, pro cultu sanctárum Imáginum 
 se Leóni Arméno, Imperatóri Iconoclástæ, constánter oppósuit; a quo 
 mulctátus exsílio, ibídem, cum per quatuórdecim annos longum duxísset 
 martyrium, migrávit ad Dóminum.
\switchcolumn
\selectlanguage{english}
In the island of Marmara, in the 
 Sea of Marmara, St. Nicephorus, bishop of Constantinople. In defence 
 of the traditions of the Fathers and of the veneration of sacred images, he 
 set himself firmly against the Iconoclast emperor Leo the Armenian, by whom 
 he was sent into exile. There he underwent a long martyrdom of 
 fourteen years and then departed for the kingdom of God.
\switchcolumn*
\selectlanguage{latin}
Romæ sancti Eugénii 
 Primi, Papæ et Confessóris.
\switchcolumn
\selectlanguage{english}
At Rome, Pope St. Eugene I, 
 Confessor.
\switchcolumn*
\selectlanguage{latin}
Trani, in Apúlia, 
 sancti Nicolái Peregríni Confessóris, cujus mirácula in Concílio Románo, cui 
 beátus Urbánus Papa Secúndus præfuit, recitáta sunt.
\switchcolumn
\selectlanguage{english}
At Trani in Apulia, St. Nicholas 
 Peregrinus, confessor, whose miracles were recounted in the Roman Council 
 under Pope Urban II.
\switchcolumn*
\selectlanguage{latin}
\end{paracol}


% ---- martyrology/mart06/mart0603.htm
\needspace{10\baselineskip}
\begin{paracol}{2}
\selectlanguage{latin}
\begin{center}{\color{gregoriocolor} Tértio Nonas Júnii. 
 Luna\dots\ }\end{center}
\switchcolumn
\selectlanguage{english}
\begin{center}{\color{gregoriocolor} The 3\textsuperscript{rd} Day of 
 June. The\dots\ Day of the Moon.}\end{center}
\end{paracol}

\noindent\begin{tabularx}{\linewidth}{*{19}{>{\centering\arraybackslash}X}}
 \textcolor{gregoriocolor}{a} & \textcolor{gregoriocolor}{b} & \textcolor{gregoriocolor}{c} & \textcolor{gregoriocolor}{d} & \textcolor{gregoriocolor}{e} & \textcolor{gregoriocolor}{f} & \textcolor{gregoriocolor}{g} & \textcolor{gregoriocolor}{h} & \textcolor{gregoriocolor}{i} & \textcolor{gregoriocolor}{k} & \textcolor{gregoriocolor}{l} & \textcolor{gregoriocolor}{m} & \textcolor{gregoriocolor}{n} & \textcolor{gregoriocolor}{p} & \textcolor{gregoriocolor}{q} & \textcolor{gregoriocolor}{r} & \textcolor{gregoriocolor}{s} & \textcolor{gregoriocolor}{t} & \textcolor{gregoriocolor}{u} \\
 7 & 8 & 9 & 10 & 11 & 12 & 13 & 14 & 15 & 16 & 17 & 18 & 19 & 20 & 21 & 22 & 23 & 24 & 25 \\
\end{tabularx}
\vspace{0.5\baselineskip}
\noindent\begin{tabularx}{\linewidth}{*{12}{>{\centering\arraybackslash}X}}
 \textcolor{gregoriocolor}{A} & \textcolor{gregoriocolor}{B} & \textcolor{gregoriocolor}{C} & \textcolor{gregoriocolor}{D} & \textcolor{gregoriocolor}{E} & F & \textcolor{gregoriocolor}{F} & \textcolor{gregoriocolor}{G} & \textcolor{gregoriocolor}{H} & \textcolor{gregoriocolor}{M} & \textcolor{gregoriocolor}{N} & \textcolor{gregoriocolor}{P} \\
 26 & 27 & 28 & 29 & 1 & 2 & 1 & 2 & 3 & 4 & 5 & 6 \\
\end{tabularx}

\begin{paracol}{2}
\selectlanguage{latin}
\lettrine[lines=2]{A}{rétii,} in Túscia, 
 sanctórum Mártyrum Pergentíni et Laurentíni fratrum, qui, in persecutióne 
 Décii, sub Tibúrtio Præside, cum essent púeri, ibídem, post dira supplícia 
 toleráta et magna mirácula osténsa, gládio cæsi sunt.
\switchcolumn
\selectlanguage{english}
\lettrine[lines=2]{A}{t} Arezzo in Tuscany, during the 
 persecution of Decius, under Governor Tiburtius, the holy martyrs 
 Pergentinus and Laurentinus, brothers, who being as yet children, were put 
 to the sword after they had endured cruel torments and performed many 
 miracles.
\switchcolumn*
\selectlanguage{latin}
Constantinópoli 
 sanctórum Mártyrum Lucilliáni et quátuor puerórum, scílicet Cláudii, Hypátii, 
 Pauli et Dionysii. His cum púeris Lucilliánus, ex idolórum sacerdóte 
 Christiánus factus, in fornácem, post vária torménta, injéctus est, sed, 
 flamma imbre exstíncta, omnes illæsi evasérunt; dénique, ipse cruci affíxus, 
 púeri autem gládio obtruncáti, sub Silváno Præside, consummáti sunt.
\switchcolumn
\selectlanguage{english}
At Constantinople, the holy martyrs 
 Lucillian and four boys, Claudius, Hypatius, Paul, and Denis. 
 Lucillian, formerly a pagan priest, but now a Christian, was cast with them 
 into a furnace after undergoing many torments, but the flames were 
 extinguished by the rain and all escaped injury. Finally their lives 
 were ended under the governor Silvanus, Lucillian by crucifixion, the 
 children by beheading.
\switchcolumn*
\selectlanguage{latin}
Córdubæ, in Hispánia, 
 beáti Isaac Mónachi, qui, pro Christi fide, gládio necátus est.
\switchcolumn
\selectlanguage{english}
At Cordova in Spain, blessed Isaac, 
 a monk who was slain by the sword for the faith of Christ.
\switchcolumn*
\selectlanguage{latin}
Constantinópoli sanctæ 
 Paulæ, Vírginis et Mártyris, quæ, cum prædictórum Mártyrum Lucilliáni ac 
 Sociórum sánguinem collígeret, ídeo, comprehénsa, virgis percússa, et in 
 ignem conjécta sed liberáta, demum et ipsa eódem loco, ubi sanctus 
 Lucilliánus crucifíxus fúerat, decolláta est.
\switchcolumn
\selectlanguage{english}
At Constantinople, St. Paula, virgin 
 and martyr, who was arrested while gathering the blood of the martyrs just 
 mentioned. She was beaten with rods and thrown into the fire, but was 
 delivered from it. She was at length beheaded in the same place where 
 St. Lucillian had been crucified.
\switchcolumn*
\selectlanguage{latin}
Carthágine sancti 
 Cæcílii Presbyteri, qui sanctum Cypriánum ad Christi fidem perdúxit.
\switchcolumn
\selectlanguage{english}
At Carthage, St. Caecilius, the 
 priest who converted St. Cyprian to the faith of Christ.
\switchcolumn*
\selectlanguage{latin}
In território 
 Aurelianénsi sancti Liphárdi, Presbyteri et Confessóris.
\switchcolumn
\selectlanguage{english}
In the diocese of Orleans, St. 
 Lifard, priest and confessor.
\switchcolumn*
\selectlanguage{latin}
Lucæ, in Túscia, sancti 
 Davíni Confessóris.
\switchcolumn
\selectlanguage{english}
At Lucca in Tuscany, St. Davinus, 
 confessor.
\switchcolumn*
\selectlanguage{latin}
Anágniæ sanctæ Olívæ 
 Vírginis.
\switchcolumn
\selectlanguage{english}
At Anagni, St. Olive, virgin.
\switchcolumn*
\selectlanguage{latin}
Lutétiæ Parisiórum 
 sanctæ Clotíldis Regínæ, cujus précibus vir ejus Clodovéus, Rex Francórum, 
 Christi fidem suscépit.
\switchcolumn
\selectlanguage{english}
At Paris, St. Clotilde, queen, by 
 whose prayers her husband, King Clovis, was converted to the faith of 
 Christ.
\switchcolumn*
\selectlanguage{latin}
\end{paracol}


% ---- martyrology/mart06/mart0604.htm
\needspace{10\baselineskip}
\begin{paracol}{2}
\selectlanguage{latin}
\begin{center}{\color{gregoriocolor} Prídie Nonas Júnii. 
 Luna\dots\ }\end{center}
\switchcolumn
\selectlanguage{english}
\begin{center}{\color{gregoriocolor} The 4\textsuperscript{th} Day of 
 June. The\dots\ Day of the Moon.}\end{center}
\end{paracol}

\noindent\begin{tabularx}{\linewidth}{*{19}{>{\centering\arraybackslash}X}}
 \textcolor{gregoriocolor}{a} & \textcolor{gregoriocolor}{b} & \textcolor{gregoriocolor}{c} & \textcolor{gregoriocolor}{d} & \textcolor{gregoriocolor}{e} & \textcolor{gregoriocolor}{f} & \textcolor{gregoriocolor}{g} & \textcolor{gregoriocolor}{h} & \textcolor{gregoriocolor}{i} & \textcolor{gregoriocolor}{k} & \textcolor{gregoriocolor}{l} & \textcolor{gregoriocolor}{m} & \textcolor{gregoriocolor}{n} & \textcolor{gregoriocolor}{p} & \textcolor{gregoriocolor}{q} & \textcolor{gregoriocolor}{r} & \textcolor{gregoriocolor}{s} & \textcolor{gregoriocolor}{t} & \textcolor{gregoriocolor}{u} \\
 8 & 9 & 10 & 11 & 12 & 13 & 14 & 15 & 16 & 17 & 18 & 19 & 20 & 21 & 22 & 23 & 24 & 25 & 26 \\
\end{tabularx}
\vspace{0.5\baselineskip}
\noindent\begin{tabularx}{\linewidth}{*{12}{>{\centering\arraybackslash}X}}
 \textcolor{gregoriocolor}{A} & \textcolor{gregoriocolor}{B} & \textcolor{gregoriocolor}{C} & \textcolor{gregoriocolor}{D} & \textcolor{gregoriocolor}{E} & F & \textcolor{gregoriocolor}{F} & \textcolor{gregoriocolor}{G} & \textcolor{gregoriocolor}{H} & \textcolor{gregoriocolor}{M} & \textcolor{gregoriocolor}{N} & \textcolor{gregoriocolor}{P} \\
 27 & 28 & 29 & 1 & 2 & 3 & 2 & 3 & 4 & 5 & 6 & 7 \\
\end{tabularx}

\begin{paracol}{2}
\selectlanguage{latin}
\lettrine[lines=2]{A}{gnóni,} in Aprútio 
 citerióre, sancti Francísci, ex nóbili Neapolitána família Carácciolo, 
 Confessóris, Congregatiónis Clericórum Regulárium Minórum Fundatóris, qui 
 mira in Deum et próximum caritáte et ardentíssimo sacræ Eucharístiæ cultus 
 propagándi stúdio flagrávit; atque a Pio Séptimo, Pontífice Máximo, 
 Sanctórum cánoni adscríptus est. Ipsíus corpus Neápolim, in Campánia, 
 translátum fuit, ibíque religiosíssime cólitur.
\switchcolumn
\selectlanguage{english}
\lettrine[lines=2]{A}{t} Agnone in Abruzzo, St. Francis of 
 the noble Neapolitan family Caracciolo, confessor, and founder of the 
 Congregation of Minor Clerks Regular. He burned with an admirable love 
 of God and of neighbour, and a most ardent desire to spread devotion to the 
 Most Holy Eucharist. His body was taken to Naples in Campania, where 
 it is religiously honoured. He was inscribed in the catalogue of the 
 saints by Pius VII.
\switchcolumn*
\selectlanguage{latin}
Romæ sanctórum Mártyrum 
 Arétii et Daciáni.
\switchcolumn
\selectlanguage{english}
At Rome, the holy martyrs Aretius 
 and Dacian.
\switchcolumn*
\selectlanguage{latin}
Sísciæ, in Illyrico, 
 sancti Quiríni Epíscopi, qui, sub Galério Præside, pro fide Christi (ut 
 Prudéntius scribit), molári saxo ad collum ligáto, in flumen præcipitátus 
 est; sed, lápide supernatánte, cum circumstántes Christiános, ne ejus 
 terreréntur supplício neve titubárent in fide, diu fuísset hortátus, ipse, 
 ut martyrii glóriam assequerétur, précibus a Deo, ut mergerétur, obtínuit.
\switchcolumn
\selectlanguage{english}
At Sissek in Illyria, in the time of 
 Governor Galerius, St. Quirinus, bishop. Prudentius relates that for 
 the faith of Christ he was thrown into a river with a millstone tied to his 
 neck. But the stone floated, and he for a long time exhorted the 
 Christians who were present not to be terrified by his punishment, nor to 
 waver in the faith, and then obtained of God by his prayers that he should 
 be drowned in order to attain the glory of martyrdom.
\switchcolumn*
\selectlanguage{latin}
Medioláni, sancti 
 Clatéi, Epíscopi Brixiénsis et Mártyris, qui, sub Neróne Imperatóre, jussu 
 Præfécti urbis illíus tentus, et, cum renuntiáre Christo nollet, multis 
 verbéribus afflíctus et cápite obtruncátus est.
\switchcolumn
\selectlanguage{english}
At Milan, in the reign of Emperor 
 Nero, St. Clateus, bishop of Brescia and martyr. By order of the 
 prefect of the city he was arrested, and when he would not deny Christ he 
 was cruelly scourged and beheaded.
\switchcolumn*
\selectlanguage{latin}
In Pannónia sanctórum 
 Mártyrum Rútili et Sociórum.
\switchcolumn
\selectlanguage{english}
In Hungary, the holy martyrs Rutilus 
 and his companions.
\switchcolumn*
\selectlanguage{latin}
Tíbure sancti Quiríni 
 Mártyris.
\switchcolumn
\selectlanguage{english}
At Tivoli, St Quirinus, martyr.
\switchcolumn*
\selectlanguage{latin}
Atrébati, in Gálliis, 
 sanctæ Saturnínæ, Vírginis et Mártyris.
\switchcolumn
\selectlanguage{english}
At Arras in France, St. Saturnina, 
 virgin and martyr.
\switchcolumn*
\selectlanguage{latin}
Constantinópoli sancti 
 Metróphanis, Epíscopi et Confessóris insígnis.
\switchcolumn
\selectlanguage{english}
At Constantinople, St. Metrophanes, 
 bishop and renowned confessor.
\switchcolumn*
\selectlanguage{latin}
Milévi, in Numídia, 
 sancti Optáti Epíscopi, doctrína et sanctitáte conspícui, quem sancti 
 Ecclésiæ Patres Augustínus et Fulgéntius suis láudibus celebrárunt.
\switchcolumn
\selectlanguage{english}
At Milevi in Numidia, St. Optatus, 
 bishop, celebrated for his learning and holiness. The holy Fathers of 
 the Church, Augustine and Fulgentius, prasied him highly.
\switchcolumn*
\selectlanguage{latin}
Verónæ sancti 
 Alexándri Epíscopi.
\switchcolumn
\selectlanguage{english}
At Verona, St. Alexander, bishop.
\switchcolumn*
\selectlanguage{latin}
\end{paracol}


% ---- martyrology/mart06/mart0605.htm
\needspace{10\baselineskip}
\begin{paracol}{2}
\selectlanguage{latin}
\begin{center}{\color{gregoriocolor} Nonis Júnii. 
 Luna\dots\ }\end{center}
\switchcolumn
\selectlanguage{english}
\begin{center}{\color{gregoriocolor} The 5\textsuperscript{th} Day of 
 June. The\dots\ Day of the Moon.}\end{center}
\end{paracol}

\noindent\begin{tabularx}{\linewidth}{*{19}{>{\centering\arraybackslash}X}}
 \textcolor{gregoriocolor}{a} & \textcolor{gregoriocolor}{b} & \textcolor{gregoriocolor}{c} & \textcolor{gregoriocolor}{d} & \textcolor{gregoriocolor}{e} & \textcolor{gregoriocolor}{f} & \textcolor{gregoriocolor}{g} & \textcolor{gregoriocolor}{h} & \textcolor{gregoriocolor}{i} & \textcolor{gregoriocolor}{k} & \textcolor{gregoriocolor}{l} & \textcolor{gregoriocolor}{m} & \textcolor{gregoriocolor}{n} & \textcolor{gregoriocolor}{p} & \textcolor{gregoriocolor}{q} & \textcolor{gregoriocolor}{r} & \textcolor{gregoriocolor}{s} & \textcolor{gregoriocolor}{t} & \textcolor{gregoriocolor}{u} \\
 9 & 10 & 11 & 12 & 13 & 14 & 15 & 16 & 17 & 18 & 19 & 20 & 21 & 22 & 23 & 24 & 25 & 26 & 27 \\
\end{tabularx}
\vspace{0.5\baselineskip}
\noindent\begin{tabularx}{\linewidth}{*{12}{>{\centering\arraybackslash}X}}
 \textcolor{gregoriocolor}{A} & \textcolor{gregoriocolor}{B} & \textcolor{gregoriocolor}{C} & \textcolor{gregoriocolor}{D} & \textcolor{gregoriocolor}{E} & F & \textcolor{gregoriocolor}{F} & \textcolor{gregoriocolor}{G} & \textcolor{gregoriocolor}{H} & \textcolor{gregoriocolor}{M} & \textcolor{gregoriocolor}{N} & \textcolor{gregoriocolor}{P} \\
 28 & 29 & 1 & 2 & 3 & 4 & 3 & 4 & 5 & 6 & 7 & 8 \\
\end{tabularx}

\begin{paracol}{2}
\selectlanguage{latin}
\lettrine[lines=2]{I}{n} Frísia sancti 
 Bonifátii, Epíscopi Moguntíni et Mártyris. Hic de Anglia Romam venit, 
 índeque a beáto Gregório Papa Secúndo in Germániam missus est ut Christi 
 fidem illis géntibus evangelizáret, et, cum máximum ibi multitúdinem, 
 præsértim Frísonum, Christiánæ religióni subjugásset, Germanórum Apóstolus 
 méruit appellári; novíssime in Frísia, a furéntibus Gentílibus gládio 
 perémptus, una cum Eóbano Coepíscopo et quibúsdam áliis servis Dei, martyrium 
 consummávit.
\switchcolumn
\selectlanguage{english}
\lettrine[lines=2]{I}{n} Friesland, St. Boniface, bishop 
 of Mainz, and martyr. He went from England to Rome, and was then sent 
 by Pope Gregory II to Germany to preach the faith of Christ to the people of 
 that country. After converting large multitudes to the Christian 
 religion, especially in Friesland, he merited the title Apostle of the 
 Germans. His martyrdom was fulfilled by being put to the sword by the 
 furious heathens, along with his fellow bishop Eobanus and some other 
 servants of God.
\switchcolumn*
\selectlanguage{latin}
Tyri, in Phœnícia, 
 sancti Doróthei Presbyteri, qui, sub Diocletiáno, multa passus est; et, 
 usque ad Juliáni témpora supérstes, sub eo, annum agens séptimum supra 
 centésimum, venerándam senéctam martyrio honestávit.
\switchcolumn
\selectlanguage{english}
At Tyre, St. Dorotheus, a priest, 
 who suffered greatly under Diocletian, but survived until the reign of 
 Julian, under whom his venerable age of one hundred and seven years was 
 crowned with martyrdom.
\switchcolumn*
\selectlanguage{latin}
In Ægypto natális 
 sanctórum Mártyrum Marciáni, Nicánoris, Apollónii et aliórum, qui, in 
 persecutióne Galérii Maximiáni, illústre martyrium consummárunt.
\switchcolumn
\selectlanguage{english}
In Egypt, the birthday of the holy 
 martyrs Marcian, Nicanor, Apollonius, and others, who suffered a glorious 
 martyrdom.
\switchcolumn*
\selectlanguage{latin}
Perúsiæ sanctórum 
 Mártyrum Floréntii, Juliáni, Cyríaci, Marcellíni et Faustíni, qui omnes, in 
 persecutióne Décii Imperatóris, cápite cæsi sunt.
\switchcolumn
\selectlanguage{english}
At Perugia, the holy martyrs 
 Florentius, Julian, Cyriacus, Marcellinus, and Faustinus, who were beheaded 
 in the persecution of Decius.
\switchcolumn*
\selectlanguage{latin}
Córdubæ, in Hispánia, 
 beáti Sáncii adolescéntis, qui, etsi in aula régia educátus, pro Christi 
 tamen fide, in persecutióne Arábica, martyrium subíre non dubitávit.
\switchcolumn
\selectlanguage{english}
At Cordova in Spain, blessed Sancho, 
 a youth brought up in the royal court, who did not hesitate to undergo 
 martyrdom for the faith of Christ during the persecution by the Arabs.
\switchcolumn*
\selectlanguage{latin}
Cæsaréæ, in Palæstína, 
 pássio sanctárum Zenáidis, Cyriæ, Valériæ et Márciæ; quæ, per multa torménta, 
 gaudéntes ad martyrium pervenérunt.
\switchcolumn
\selectlanguage{english}
At Caesarea in Palestine, the 
 martyrdom of the Saints Zenaides, Cyria, Valeria, and Marcia, who joyfully 
 attained martyrdom through many torments.
\switchcolumn*
\selectlanguage{latin}
\end{paracol}


% ---- martyrology/mart06/mart0606.htm
\needspace{10\baselineskip}
\begin{paracol}{2}
\selectlanguage{latin}
\begin{center}{\color{gregoriocolor} Octávo Idus Júnii. 
 Luna\dots\ }\end{center}
\switchcolumn
\selectlanguage{english}
\begin{center}{\color{gregoriocolor} The 6\textsuperscript{th} Day of 
 June. The\dots\ Day of the Moon.}\end{center}
\end{paracol}

\noindent\begin{tabularx}{\linewidth}{*{19}{>{\centering\arraybackslash}X}}
 \textcolor{gregoriocolor}{a} & \textcolor{gregoriocolor}{b} & \textcolor{gregoriocolor}{c} & \textcolor{gregoriocolor}{d} & \textcolor{gregoriocolor}{e} & \textcolor{gregoriocolor}{f} & \textcolor{gregoriocolor}{g} & \textcolor{gregoriocolor}{h} & \textcolor{gregoriocolor}{i} & \textcolor{gregoriocolor}{k} & \textcolor{gregoriocolor}{l} & \textcolor{gregoriocolor}{m} & \textcolor{gregoriocolor}{n} & \textcolor{gregoriocolor}{p} & \textcolor{gregoriocolor}{q} & \textcolor{gregoriocolor}{r} & \textcolor{gregoriocolor}{s} & \textcolor{gregoriocolor}{t} & \textcolor{gregoriocolor}{u} \\
 10 & 11 & 12 & 13 & 14 & 15 & 16 & 17 & 18 & 19 & 20 & 21 & 22 & 23 & 24 & 25 & 26 & 27 & 28 \\
\end{tabularx}
\vspace{0.5\baselineskip}
\noindent\begin{tabularx}{\linewidth}{*{12}{>{\centering\arraybackslash}X}}
 \textcolor{gregoriocolor}{A} & \textcolor{gregoriocolor}{B} & \textcolor{gregoriocolor}{C} & \textcolor{gregoriocolor}{D} & \textcolor{gregoriocolor}{E} & F & \textcolor{gregoriocolor}{F} & \textcolor{gregoriocolor}{G} & \textcolor{gregoriocolor}{H} & \textcolor{gregoriocolor}{M} & \textcolor{gregoriocolor}{N} & \textcolor{gregoriocolor}{P} \\
 29 & 1 & 2 & 3 & 4 & 5 & 4 & 5 & 6 & 7 & 8 & 9 \\
\end{tabularx}

\begin{paracol}{2}
\selectlanguage{latin}
\lettrine[lines=2]{M}{agdebúrgi} sancti 
 Norbérti, ejúsdem civitátis Epíscopi et Confessóris, qui Fundátor exstitit 
 Ordinis Præmonstraténsis.
\switchcolumn
\selectlanguage{english}
\lettrine[lines=2]{A}{t} Magdeburg, St. Norbert, bishop of 
 that city, confessor. He was the founder of the Premonstratensian 
 Order.
\switchcolumn*
\selectlanguage{latin}
Cæsaréæ, in Palæstína, 
 natális beáti Philíppi, qui fuit unus de septem primis Diáconis. Hic, 
 signis et prodígiis clarus, Samaríam ad Christi fidem convértit, et Regínæ 
 Æthíopum Cándacis Eunúchum baptizávit, ac demum apud Cæsaréam requiévit. 
 Juxta ipsum tres Vírgines, ejus fíliæ ac Prophetíssæ, tumulátæ jacent; nam 
 quarta fília ejus, plena Spíritu Sancto, Ephesi occúbuit.
\switchcolumn
\selectlanguage{english}
At Caesarea in Palestine, the 
 birthday of blessed Philip, one of the first seven deacons. He was 
 renowned for miracles and prodigies. He converted Samaria to the faith 
 of Christ, baptized the eunuch of Candace, queen of Ethiopia, and finally 
 rested in peace at Caesarea. Near him are buried three of his 
 daughters, virgins and prophetesses. His fourth daughter died at 
 Ephesus, filled with the Holy Ghost.
\switchcolumn*
\selectlanguage{latin}
Romæ sancti Artémii, 
 cum uxóre sua Cándida, et fília Paulína. Ex his Artémius, ad 
 prædicatiónem et mirácula sancti Petri Exorcístæ, ad Christum convérsus, et, 
 cum omni domo sua, a sancto Marcellíno Presbytero baptizátus, Seréni Judicis 
 jussu plumbátis cæsus et gládio percússus est; uxor vero ejus et fília, in 
 cryptam impúlsæ, lapídibus ruderibúsque sunt óbrutæ.
\switchcolumn
\selectlanguage{english}
At Rome, St. Artemius, with his wife 
 Candida and his daughter Paulina. Artemius became a believer through 
 the preaching and miracles of St. Peter the Exorcist, who was baptized with 
 all his household by the priest St. Marcellinus. By order of Judge 
 Serenus, he was scourged with leaded whips, and then slain with the sword. 
 His wife and daughter were forced into a pit and covered with stones and 
 earth.
\switchcolumn*
\selectlanguage{latin}
In agro Bononiénsi 
 sancti Alexándri, Epíscopi Fæsuláni et Mártyris; qui, rédiens ex urbe Papía, 
 ubi Ecclésiæ suæ bona apud Longobardórum Regem ab usurpatóribus vindicáverat, 
 ab his in Rhenum flúvium est dejéctus et aquis præfocátus.
\switchcolumn
\selectlanguage{english}
In the district of Bologna, St. 
 Alexander, bishop of Fiesole and martyr. While returning from the town 
 of Pavia, where he had defended the title to the goods of his church before 
 the Lombard king against those taking them away, he was seized by the 
 usurpers, cast into the Rhine river, and drowned.
\switchcolumn*
\selectlanguage{latin}
Tarsi, in Cilícia, 
 sanctórum Mártyrum vigínti, qui, tempóribus Diocletiáni et Maximiáni, sub 
 Simplício Júdice, per divérsa torménta glorificavérunt Deum in corpóribus 
 suis.
\switchcolumn
\selectlanguage{english}
At Tarsus in Cilicia, in the time of 
 Emperors Diocletian and Maximian, and the governor Simplicius, twenty holy 
 martyrs, who, through various torments to their bodies, glorified God.
\switchcolumn*
\selectlanguage{latin}
Noviodúni, in Gálliis, 
 sanctórum Mártyrum Amántii, Alexándri et Sociórum.
\switchcolumn
\selectlanguage{english}
At Noyon in France, the holy martyrs 
 Amantius, Alexander, and their companions.
\switchcolumn*
\selectlanguage{latin}
Medioláni deposítio 
 sancti Eustórgii Secúndi, Epíscopi et Confessóris.
\switchcolumn
\selectlanguage{english}
At Milan, the death of St. 
 Eustorgius II, bishop and confessor.
\switchcolumn*
\selectlanguage{latin}
Verónæ sancti Joánnis 
 Epíscopi.
\switchcolumn
\selectlanguage{english}
At Verona, the bishop St. John.
\switchcolumn*
\selectlanguage{latin}
Vesontióne, in Gálliis, 
 sancti Cláudii Epíscopi.
\switchcolumn
\selectlanguage{english}
At Besançon, France, the bishop St. 
 Claudius.
\switchcolumn*
\selectlanguage{latin}
\end{paracol}


% ---- martyrology/mart06/mart0607.htm
\needspace{10\baselineskip}
\begin{paracol}{2}
\selectlanguage{latin}
\begin{center}{\color{gregoriocolor} Séptimo Idus Júnii. 
 Luna\dots\ }\end{center}
\switchcolumn
\selectlanguage{english}
\begin{center}{\color{gregoriocolor} The 7\textsuperscript{th} Day of 
 June. The\dots\ Day of the Moon.}\end{center}
\end{paracol}

\noindent\begin{tabularx}{\linewidth}{*{19}{>{\centering\arraybackslash}X}}
 \textcolor{gregoriocolor}{a} & \textcolor{gregoriocolor}{b} & \textcolor{gregoriocolor}{c} & \textcolor{gregoriocolor}{d} & \textcolor{gregoriocolor}{e} & \textcolor{gregoriocolor}{f} & \textcolor{gregoriocolor}{g} & \textcolor{gregoriocolor}{h} & \textcolor{gregoriocolor}{i} & \textcolor{gregoriocolor}{k} & \textcolor{gregoriocolor}{l} & \textcolor{gregoriocolor}{m} & \textcolor{gregoriocolor}{n} & \textcolor{gregoriocolor}{p} & \textcolor{gregoriocolor}{q} & \textcolor{gregoriocolor}{r} & \textcolor{gregoriocolor}{s} & \textcolor{gregoriocolor}{t} & \textcolor{gregoriocolor}{u} \\
 11 & 12 & 13 & 14 & 15 & 16 & 17 & 18 & 19 & 20 & 21 & 22 & 23 & 24 & 25 & 26 & 27 & 28 & 29 \\
\end{tabularx}
\vspace{0.5\baselineskip}
\noindent\begin{tabularx}{\linewidth}{*{12}{>{\centering\arraybackslash}X}}
 \textcolor{gregoriocolor}{A} & \textcolor{gregoriocolor}{B} & \textcolor{gregoriocolor}{C} & \textcolor{gregoriocolor}{D} & \textcolor{gregoriocolor}{E} & F & \textcolor{gregoriocolor}{F} & \textcolor{gregoriocolor}{G} & \textcolor{gregoriocolor}{H} & \textcolor{gregoriocolor}{M} & \textcolor{gregoriocolor}{N} & \textcolor{gregoriocolor}{P} \\
 1 & 2 & 3 & 4 & 5 & 6 & 5 & 6 & 7 & 8 & 9 & 10 \\
\end{tabularx}

\begin{paracol}{2}
\selectlanguage{latin}
\lettrine[lines=2]{C}{onstantinópoli} natális 
 sancti Pauli, ejúsdem civitátis Epíscopi, qui sæpe ab Ariánis ob fidem 
 cathólicam pulsus, et a sancto Júlio Primo, Románo Pontífice, restitútus est; 
 ac demum, ab Ariáno Imperatóre Constántio relegátus ad Cucúsum, Cappadóciæ 
 oppídulum, ibídem, Arianórum insídiis crudéliter strangulátus, ad cæléstia 
 regna migrávit. Ipsíus autem corpus, Theodósio Imperatóre, 
 Constantinópolim summo honóre, translátum fuit.
\switchcolumn
\selectlanguage{english}
\lettrine[lines=2]{A}{t} Constantinople, the birthday of 
 St. Paul, bishop of that city. For the Catholic faith, he was often 
 driven out of his see by the Arians, but restored to it by the Roman 
 Pontiff, St. Julius I. Finally the Arian emperor Constantius banished 
 him to Cucusum, a small town of Cappadocia. There, by the intrigue of 
 the Arians, he was barbarously strangled, and thus departed for the heavenly 
 kingdom. His body was taken to Constantinople with the greatest honour 
 during the reign of Emperor Theodosius.
\switchcolumn*
\selectlanguage{latin}
Córdubæ, in Hispánia, 
 sanctórum Monachórum et Mártyrum Petri Presbyteri, Wallabónsi Diáconi, 
 Sabiniáni, Wistremúndi, Habéntii et Jeremíæ, qui pro Christo, in 
 persecutióne Arábica, sunt juguláti.
\switchcolumn
\selectlanguage{english}
At Cordova in Spain, the holy 
 martyrs Peter, a priest, Wallabonsus, a deacon, Sabinianus, Wistremund, 
 Habentius, and Jeremias, all of whom were monks. Their throats were 
 cut at the time of the Arab persecution because they had confessed Christ.
\switchcolumn*
\selectlanguage{latin}
Hermópoli, in Ægypto, 
 sancti Lycariónis Mártyris, qui laniátus, virgis férreis ignítis cæsus, 
 áliaque sævíssima passus est, ac demum, gládio percússus, martyrium 
 consummávit.
\switchcolumn
\selectlanguage{english}
At Hermopolis in Egypt, St. Licarion, 
 martyr, who had his body lacerated, was scourged with heated iron rods, and 
 endured other horrible torments, after which his martyrdom was completed by 
 beheading.
\switchcolumn*
\selectlanguage{latin}
Placéntiæ sancti 
 Antónii Maríæ Gianélli, Bobiénsis Epíscopi, Fundatóris Congregatiónis 
 Filiárum Maríæ sanctíssimæ ab Horto nuncupatárum, quem Pius Papa Duodécimus 
 inter sanctos Cælites adnumerávit.
\switchcolumn
\selectlanguage{english}
At Placentia, St. Anthony Mary 
 Gianelli, bishop of Bobbio, and founder of the Congregation of Sisters of 
 our Lady of the Garden. Pope Pius XII numbered him among the saints of 
 heaven.
\switchcolumn*
\selectlanguage{latin}
In Anglia sancti 
 Robérti Abbátis, ex Ordine Cisterciénsi.
\switchcolumn
\selectlanguage{english}
In England, St. Robert, an abbot of 
 the Cistercian Order.
\switchcolumn*
\selectlanguage{latin}
\end{paracol}


% ---- martyrology/mart06/mart0608.htm
\needspace{10\baselineskip}
\begin{paracol}{2}
\selectlanguage{latin}
\begin{center}{\color{gregoriocolor} Sexto Idus Júnii. 
 Luna\dots\ }\end{center}
\switchcolumn
\selectlanguage{english}
\begin{center}{\color{gregoriocolor} The 8\textsuperscript{th} Day of 
 June. The\dots\ Day of the Moon.}\end{center}
\end{paracol}

\noindent\begin{tabularx}{\linewidth}{*{19}{>{\centering\arraybackslash}X}}
 \textcolor{gregoriocolor}{a} & \textcolor{gregoriocolor}{b} & \textcolor{gregoriocolor}{c} & \textcolor{gregoriocolor}{d} & \textcolor{gregoriocolor}{e} & \textcolor{gregoriocolor}{f} & \textcolor{gregoriocolor}{g} & \textcolor{gregoriocolor}{h} & \textcolor{gregoriocolor}{i} & \textcolor{gregoriocolor}{k} & \textcolor{gregoriocolor}{l} & \textcolor{gregoriocolor}{m} & \textcolor{gregoriocolor}{n} & \textcolor{gregoriocolor}{p} & \textcolor{gregoriocolor}{q} & \textcolor{gregoriocolor}{r} & \textcolor{gregoriocolor}{s} & \textcolor{gregoriocolor}{t} & \textcolor{gregoriocolor}{u} \\
 12 & 13 & 14 & 15 & 16 & 17 & 18 & 19 & 20 & 21 & 22 & 23 & 24 & 25 & 26 & 27 & 28 & 29 & 1 \\
\end{tabularx}
\vspace{0.5\baselineskip}
\noindent\begin{tabularx}{\linewidth}{*{12}{>{\centering\arraybackslash}X}}
 \textcolor{gregoriocolor}{A} & \textcolor{gregoriocolor}{B} & \textcolor{gregoriocolor}{C} & \textcolor{gregoriocolor}{D} & \textcolor{gregoriocolor}{E} & F & \textcolor{gregoriocolor}{F} & \textcolor{gregoriocolor}{G} & \textcolor{gregoriocolor}{H} & \textcolor{gregoriocolor}{M} & \textcolor{gregoriocolor}{N} & \textcolor{gregoriocolor}{P} \\
 2 & 3 & 4 & 5 & 6 & 7 & 6 & 7 & 8 & 9 & 10 & 11 \\
\end{tabularx}

\begin{paracol}{2}
\selectlanguage{latin}
\lettrine[lines=2]{A}{quis} in Gállia, sancti 
 Maximíni, qui éxstitit primus ejúsdem civitátis Epíscopus, ac Dómini 
 discípulus fuísse tráditur.
\switchcolumn
\selectlanguage{english}
\lettrine[lines=2]{A}{t} Aix in France, St. Maximin, first 
 bishop of that city, who is said to have been a disciple of the Lord.
\switchcolumn*
\selectlanguage{latin}
Eódem die sanctæ 
 Callíopæ Mártyris, quæ, ob Christi fidem, abscíssis mammis atque adústis 
 cárnibus, super téstulas volutáta, demum, truncáta cápite, martyrii palmam 
 accépit.
\switchcolumn
\selectlanguage{english}
On the same day, St. Calliopa, 
 martyr, who for the faith of Christ received the palm of martyrdom. 
 Her breasts had been cut away, her flesh burned, she was rolled on broken 
 pottery, and was at last beheaded.
\switchcolumn*
\selectlanguage{latin}
Eboráci, in Anglia, 
 sancti Willhélmi, Epíscopi et Confessóris, qui, inter cétera ad ejus 
 sepúlcrum patráta mirácula, tres mórtuos suscitávit, atque ab Honório Papa 
 Tértio in Sanctórum cánonem relátus est.
\switchcolumn
\selectlanguage{english}
At York in England, St. William, 
 archbishop and confessor, who, among other miracles wrought at his tomb, 
 raised three persons from the dead. He was placed in the calendar of 
 the saints by Pope Honorius III.
\switchcolumn*
\selectlanguage{latin}
Apud Suessiónes, in 
 Gálliis, natális sancti Medárdi, Epíscopi Novioménsis; cujus vita et mors 
 pretiósa gloriósis miráculis commendátur.
\switchcolumn
\selectlanguage{english}
At Soissons in France, the birthday 
 of St. Medard, bishop of Noyon, whose life and precious death are commended 
 by glorious miracles.
\switchcolumn*
\selectlanguage{latin}
Rotómagi sancti 
 Gildárdi Epíscopi, qui fuit frater ejúsdem sancti Medárdi. Ambo autem 
 fratres, eódem die nati eodémque die Epíscopi consecráti, uno quoque die de 
 hac vita subtrácti, simul in cælum migrárunt.
\switchcolumn
\selectlanguage{english}
At Rouen, St. Gildard, bishop, 
 brother of this same St. Medard. They were born on the same day, 
 consecrated bishops at the same time, and were taken from this life on the 
 same day, entering heaven together.
\switchcolumn*
\selectlanguage{latin}
Apud Sénonas sancti 
 Heráclii Epíscopi.
\switchcolumn
\selectlanguage{english}
At Sens, the bishop St. Heraclius.
\switchcolumn*
\selectlanguage{latin}
Metis, in Gállia, 
 sancti Clodúlphi Epíscopi.
\switchcolumn
\selectlanguage{english}
At Metz, the bishop St. Clodulph.
\switchcolumn*
\selectlanguage{latin}
In Picéno sancti 
 Severíni, Septempedáni Epíscopi.
\switchcolumn
\selectlanguage{english}
In Piceno, St. Severin, bishop of 
 Septempeda.
\switchcolumn*
\selectlanguage{latin}
In Sardínia sancti 
 Sallustiáni Confessóris.
\switchcolumn
\selectlanguage{english}
In Sardinia, St. Sallustian, 
 confessor.
\switchcolumn*
\selectlanguage{latin}
Cameríni sancti 
 Victoríni Confessóris, qui fuit prædícti sancti Severíni, Septempedáni 
 Epíscopi, germánus frater.
\switchcolumn
\selectlanguage{english}
At Camerino, St. Victorinus, 
 confessor, the twin brother of St. Severin, bishop of Septempeda.
\switchcolumn*
\selectlanguage{latin}
\end{paracol}


% ---- martyrology/mart06/mart0609.htm
\needspace{10\baselineskip}
\begin{paracol}{2}
\selectlanguage{latin}
\begin{center}{\color{gregoriocolor} Quinto Idus Júnii. 
 Luna\dots\ }\end{center}
\switchcolumn
\selectlanguage{english}
\begin{center}{\color{gregoriocolor} The 9\textsuperscript{th} Day of 
 June. The\dots\ Day of the Moon.}\end{center}
\end{paracol}

\noindent\begin{tabularx}{\linewidth}{*{19}{>{\centering\arraybackslash}X}}
 \textcolor{gregoriocolor}{a} & \textcolor{gregoriocolor}{b} & \textcolor{gregoriocolor}{c} & \textcolor{gregoriocolor}{d} & \textcolor{gregoriocolor}{e} & \textcolor{gregoriocolor}{f} & \textcolor{gregoriocolor}{g} & \textcolor{gregoriocolor}{h} & \textcolor{gregoriocolor}{i} & \textcolor{gregoriocolor}{k} & \textcolor{gregoriocolor}{l} & \textcolor{gregoriocolor}{m} & \textcolor{gregoriocolor}{n} & \textcolor{gregoriocolor}{p} & \textcolor{gregoriocolor}{q} & \textcolor{gregoriocolor}{r} & \textcolor{gregoriocolor}{s} & \textcolor{gregoriocolor}{t} & \textcolor{gregoriocolor}{u} \\
 13 & 14 & 15 & 16 & 17 & 18 & 19 & 20 & 21 & 22 & 23 & 24 & 25 & 26 & 27 & 28 & 29 & 1 & 2 \\
\end{tabularx}
\vspace{0.5\baselineskip}
\noindent\begin{tabularx}{\linewidth}{*{12}{>{\centering\arraybackslash}X}}
 \textcolor{gregoriocolor}{A} & \textcolor{gregoriocolor}{B} & \textcolor{gregoriocolor}{C} & \textcolor{gregoriocolor}{D} & \textcolor{gregoriocolor}{E} & F & \textcolor{gregoriocolor}{F} & \textcolor{gregoriocolor}{G} & \textcolor{gregoriocolor}{H} & \textcolor{gregoriocolor}{M} & \textcolor{gregoriocolor}{N} & \textcolor{gregoriocolor}{P} \\
 3 & 4 & 5 & 6 & 7 & 8 & 7 & 8 & 9 & 10 & 11 & 12 \\
\end{tabularx}

\begin{paracol}{2}
\selectlanguage{latin}
\lettrine[lines=2]{N}{oménti} in Sabínis, 
 natális sanctórum Mártyrum Primi et Feliciáni fratrum, sub Diocletiáno et 
 Maximiáno Imperatóribus. Hi gloriósi Mártyres, cum longævam in Dómino 
 vitam duxíssent, et nunc simul pária, nunc singillátim divérsa et exquisíta 
 pertulíssent torménta, ambo tandem felícis pugnæ cursum, a Nomentáno Præside 
 Promóto animadvérsi gládio, consummavérunt. Ipsórum autem Mártyrum 
 córpora, póstea Romam transláta, in Ecclésia sancti Stéphani Protomártyris, 
 in monte Cælio, honorífice collocáta sunt.
\switchcolumn
\selectlanguage{english}
\lettrine[lines=2]{A}{t} Nomento in the Sabine Hills, the 
 birthday of the holy martyrs Primus and Felician, under the emperors 
 Diocletian and Maximian. These glorious martyrs lived long in the 
 service of the Lord, and endured sometimes together, sometimes separately, 
 various cruel torments. They were finally beheaded by Promotus, 
 governor of Nomento, and thus happily ended their trial. Their bodies 
 were afterwards translated to Rome and honorably buried in the Church of St. 
 Stephen the Protomartyr on the Cælian Hill.
\switchcolumn*
\selectlanguage{latin}
Agénni in Gállia, 
 pássio sancti Vincéntii, Levítæ et Mártyris, qui, ob Christi fidem, 
 verbéribus diríssime cæsus et gládio decollátus est.
\switchcolumn
\selectlanguage{english}
At Agen in France, the passion of St. Vincent, deacon and 
 martyr. For the faith of Christ, he was cruelly scourged and then 
 beheaded.
\switchcolumn*
\selectlanguage{latin}
Apud Antiochíam sanctæ 
 Pelágiæ, Vírginis et Mártyris, quam sancti Ambrósius et Joánnes Chrysóstomus 
 magnis éfferunt láudibus.
\switchcolumn
\selectlanguage{english}
At Antioch, St. Pelagia, virgin and 
 martyr, who has been eulogized by St. Ambrose and St. John Chrysostom.
\switchcolumn*
\selectlanguage{latin}
Syracúsis, in Sicília, 
 sancti Maximiáni Epíscopi, cujus sanctus Gregórius Papa sæpius méminit.
\switchcolumn
\selectlanguage{english}
At Syracuse in Sicily, Bishop St. 
 Maximian, who is frequently mentioned by Pope St. Gregory.
\switchcolumn*
\selectlanguage{latin}
Andriæ, in Apúlia, 
 sancti Richárdi, qui fuit primus ejúsdem civitátis Epíscopus, et miráculis 
 cláruit.
\switchcolumn
\selectlanguage{english}
At Andria in Apulia, St. Richard, 
 first bishop of that city, who is famed for his miracles.
\switchcolumn*
\selectlanguage{latin}
In Ióna, Scótiæ ínsula, 
 sancti Colúmbæ, Presbyteri et Abbátis.
\switchcolumn
\selectlanguage{english}
In the island of Iona in Scotland, 
 St. Columba, priest and confessor.
\switchcolumn*
\selectlanguage{latin}
Edéssæ, in Syria, 
 sancti Juliáni Mónachi, cujus præclára gesta sanctus Ephræm Diáconus 
 scripsit.
\switchcolumn
\selectlanguage{english}
At Edessa in Syria, St. Julian, a 
 monk whose memorable deeds have been related by the deacon St. Ephraem.
\switchcolumn*
\selectlanguage{latin}
\end{paracol}


% ---- martyrology/mart06/mart0610.htm
\needspace{10\baselineskip}
\begin{paracol}{2}
\selectlanguage{latin}
\begin{center}{\color{gregoriocolor} Quarto Idus Júnii. 
 Luna\dots\ }\end{center}
\switchcolumn
\selectlanguage{english}
\begin{center}{\color{gregoriocolor} The 10\textsuperscript{th} Day of 
 June. The\dots\ Day of the Moon.}\end{center}
\end{paracol}

\noindent\begin{tabularx}{\linewidth}{*{19}{>{\centering\arraybackslash}X}}
 \textcolor{gregoriocolor}{a} & \textcolor{gregoriocolor}{b} & \textcolor{gregoriocolor}{c} & \textcolor{gregoriocolor}{d} & \textcolor{gregoriocolor}{e} & \textcolor{gregoriocolor}{f} & \textcolor{gregoriocolor}{g} & \textcolor{gregoriocolor}{h} & \textcolor{gregoriocolor}{i} & \textcolor{gregoriocolor}{k} & \textcolor{gregoriocolor}{l} & \textcolor{gregoriocolor}{m} & \textcolor{gregoriocolor}{n} & \textcolor{gregoriocolor}{p} & \textcolor{gregoriocolor}{q} & \textcolor{gregoriocolor}{r} & \textcolor{gregoriocolor}{s} & \textcolor{gregoriocolor}{t} & \textcolor{gregoriocolor}{u} \\
 14 & 15 & 16 & 17 & 18 & 19 & 20 & 21 & 22 & 23 & 24 & 25 & 26 & 27 & 28 & 29 & 1 & 2 & 3 \\
\end{tabularx}
\vspace{0.5\baselineskip}
\noindent\begin{tabularx}{\linewidth}{*{12}{>{\centering\arraybackslash}X}}
 \textcolor{gregoriocolor}{A} & \textcolor{gregoriocolor}{B} & \textcolor{gregoriocolor}{C} & \textcolor{gregoriocolor}{D} & \textcolor{gregoriocolor}{E} & F & \textcolor{gregoriocolor}{F} & \textcolor{gregoriocolor}{G} & \textcolor{gregoriocolor}{H} & \textcolor{gregoriocolor}{M} & \textcolor{gregoriocolor}{N} & \textcolor{gregoriocolor}{P} \\
 4 & 5 & 6 & 7 & 8 & 9 & 8 & 9 & 10 & 11 & 12 & 13 \\
\end{tabularx}

\begin{paracol}{2}
\selectlanguage{latin}
\lettrine[lines=2]{S}{anctæ} Margarítæ Víduæ, 
 Scotórum Regínæ, quæ sextodécimo Kaléndas Decémbris obdormívit in Dómino.
\switchcolumn
\selectlanguage{english}
\lettrine[lines=2]{S}{t.} Margaret, widow, queen of 
 Scotland, who slept in the Lord on the 16th of November.
\switchcolumn*
\selectlanguage{latin}
Romæ, via Salária, 
 pássio beáti Getúlii, claríssimi et doctíssimi viri, ac susceptórum e sancta 
 uxóre Symphorósa beatórum septem fratrum Mártyrum patris, ejúsque Sociórum 
 Cæreális, Amántii et Primitívi. Hi omnes, Hadriáni Imperatóris jussu, 
 a Licínio Consulári tenti, primum cæsi sunt, deínde in cárcerem trusi; 
 postrémum, incéndio tráditi, sed nullo modo ab igne læsi, martyrium suum, 
 fústibus illíso cápite, complevérunt. Ipsórum autem córpora Symphorósa, 
 beáti Getúlii uxor, collégit, et in arenário prædii sui honorífice sepelívit.
\switchcolumn
\selectlanguage{english}
At Rome, on the Salarian Way, the 
 martyrdom of blessed Getulius, a very learned nobleman, and his companions, 
 Caerealis, Amantius, and Primitivus. By order of Emperor Hadrian they 
 were arrested by the ex-consul Licinius, scourged, thrown into prison, and 
 then delivered to the flames. But the fire did not injure them, and 
 their heads were crushed with clubs, thus ending their martyrdom. 
 Their bodies were taken by Symphorosa, wife of blessed Getulius, and 
 reverently interred on her own estate.
\switchcolumn*
\selectlanguage{latin}
Item Romæ, via Aurélia, natális sanctórum Basílidis, Trípodis, Mándalis et aliórum vigínti Mártyrum, 
 sub Imperatóre Aureliáno et Urbis Præfécto Platóne.
\switchcolumn
\selectlanguage{english}
Also at Rome, on the Aurelian Way, 
 the birthday of the Saints Basilides, Tripos, Mandal, and twenty other 
 martyrs, under Emperor Aurelian and Plato, the governor of the city.
\switchcolumn*
\selectlanguage{latin}
Neápoli, in Campánia, 
 sancti Máximi, Epíscopi et Mártyris; qui, ob strénuam Nicǽnæ fídei 
 confessiónem, a Constántio Imperatóre in exsílium pulsus, ibídem, ærúmnis 
 conféctus, decéssit.
\switchcolumn
\selectlanguage{english}
At Naples in Campania, St. Maximus, 
 bishop and martyr. For having vigorously defended the Nicene Creed he 
 was exiled by Emperor Constantius, where he died worn out by his trials.
\switchcolumn*
\selectlanguage{latin}
Prusæ, in Bithynia, 
 sancti Timóthei, Epíscopi et Mártyris, qui, sub Juliáno Apóstata, cum 
 Christum ejuráre noluísset, idcírco, ipsíus Imperatóris jussu, cápite 
 abscíssus est.
\switchcolumn
\selectlanguage{english}
At Prusias in Bithynia, St. Timothy, 
 bishop and martyr. He was beheaded during the reign of Julian the 
 Apostate because he refused to deny Christ.
\switchcolumn*
\selectlanguage{latin}
Colóniæ Agrippínæ 
 sancti Mauríni, Abbátis et Mártyris.
\switchcolumn
\selectlanguage{english}
At Cologne, St. Maurinus, abbot and martyr
\switchcolumn*
\selectlanguage{latin}
Nicomedíæ sancti 
 Zacharíæ Mártyris.
\switchcolumn
\selectlanguage{english}
At Nicomedia, the martyr St. 
 Zachary.
\switchcolumn*
\selectlanguage{latin}
In Hispánia sanctórum 
 Mártyrum Críspuli et Restitúti.
\switchcolumn
\selectlanguage{english}
In Spain, the holy martyrs Crispulus 
 and Restitutus.
\switchcolumn*
\selectlanguage{latin}
In Africa sanctórum 
 Mártyrum Arésii, Rogáti et aliórum quíndecim.
\switchcolumn
\selectlanguage{english}
In Africa, the holy martyrs Aresius, 
 Rogatus, and fifteen others.
\switchcolumn*
\selectlanguage{latin}
Petræ, in Arábia, 
 sancti Astérii Epíscopi, qui, ob fidem cathólicam, ab Ariánis multa passus 
 et ab Imperatóre Constántio in Africam relegátus est, ac tandem, in 
 Ecclésiam suam restitútus, Conféssor gloriósus occúbuit.
\switchcolumn
\selectlanguage{english}
At Petra in Africa, St. Asterius, a 
 bishop who suffered greatly for the Catholic faith at the hands of the 
 Arians. He was banished to Africa by Emperor Constantius, and there 
 died as a glorious confessor.
\switchcolumn*
\selectlanguage{latin}
Antisiodóri sancti 
 Censúrii Epíscopi.
\switchcolumn
\selectlanguage{english}
At Auxerre, St. Censurius, bishop.
\switchcolumn*
\selectlanguage{latin}
\end{paracol}


% ---- martyrology/mart06/mart0611.htm
\needspace{10\baselineskip}
\begin{paracol}{2}
\selectlanguage{latin}
\begin{center}{\color{gregoriocolor} Tértio Idus Júnii. 
 Luna\dots\ }\end{center}
\switchcolumn
\selectlanguage{english}
\begin{center}{\color{gregoriocolor} The 11\textsuperscript{th} Day of 
 June. The\dots\ Day of the Moon.}\end{center}
\end{paracol}

\noindent\begin{tabularx}{\linewidth}{*{19}{>{\centering\arraybackslash}X}}
 \textcolor{gregoriocolor}{a} & \textcolor{gregoriocolor}{b} & \textcolor{gregoriocolor}{c} & \textcolor{gregoriocolor}{d} & \textcolor{gregoriocolor}{e} & \textcolor{gregoriocolor}{f} & \textcolor{gregoriocolor}{g} & \textcolor{gregoriocolor}{h} & \textcolor{gregoriocolor}{i} & \textcolor{gregoriocolor}{k} & \textcolor{gregoriocolor}{l} & \textcolor{gregoriocolor}{m} & \textcolor{gregoriocolor}{n} & \textcolor{gregoriocolor}{p} & \textcolor{gregoriocolor}{q} & \textcolor{gregoriocolor}{r} & \textcolor{gregoriocolor}{s} & \textcolor{gregoriocolor}{t} & \textcolor{gregoriocolor}{u} \\
 15 & 16 & 17 & 18 & 19 & 20 & 21 & 22 & 23 & 24 & 25 & 26 & 27 & 28 & 29 & 1 & 2 & 3 & 4 \\
\end{tabularx}
\vspace{0.5\baselineskip}
\noindent\begin{tabularx}{\linewidth}{*{12}{>{\centering\arraybackslash}X}}
 \textcolor{gregoriocolor}{A} & \textcolor{gregoriocolor}{B} & \textcolor{gregoriocolor}{C} & \textcolor{gregoriocolor}{D} & \textcolor{gregoriocolor}{E} & F & \textcolor{gregoriocolor}{F} & \textcolor{gregoriocolor}{G} & \textcolor{gregoriocolor}{H} & \textcolor{gregoriocolor}{M} & \textcolor{gregoriocolor}{N} & \textcolor{gregoriocolor}{P} \\
 5 & 6 & 7 & 8 & 9 & 10 & 9 & 10 & 11 & 12 & 13 & 14 \\
\end{tabularx}

\begin{paracol}{2}
\selectlanguage{latin}
\lettrine[lines=2]{S}{alamínæ,} in Cypro, 
 natális sancti Bárnabæ Apóstoli, qui, natióne Cyprius, cum Paulo Géntium 
 Apóstolus a discípulis ordinátus, multas regiónes cum eo peragrávit, 
 injúnctum sibi opus Evangélicæ prædicatiónis exércens; postrémo, Cyprum 
 proféctus, ibi Apostolátum suum glorióso martyrio decorávit. Ejus 
 corpus, témpore Zenónis Imperatóris, ipso Bárnaba revelánte, repértum est, 
 una cum códice Evangélii sancti Matthǽi, ejúsdem Bárnabæ manu descrípto.
\switchcolumn
\selectlanguage{english}
\lettrine[lines=2]{A}{t} Salamina in Cyprus, the birthday 
 of the apostle St. Barnabas, a native of that place. He was ordained 
 by the disciples as Apostle of the Gentiles with St. Paul, and travelled 
 with him over many regions, exercising the function committed unto him of 
 preaching the Gospel. At last he went back to Cyprus, where he 
 ennobled his apostolate by a glorious martyrdom. His body was found by 
 his own revelation, in the time of Emperor Zeno, together with a copy of St. 
 Matthew's Gospel written with his own hand.
\switchcolumn*
\selectlanguage{latin}
Salmánticæ, in 
 Hispánia, item natális sancti Joánnis a sancto Facúndo, ex Eremitárum sancti 
 Augustíni Ordine, Confessóris; qui fídei zelo, sanctimónia vitæ ac miráculis 
 cláruit. Ipsíus autem festívitas sequénti die celebrátur.
\switchcolumn
\selectlanguage{english}
At Salamanca in Spain, St. John of St. Facundus, a confessor of the Order of the Hermits of St. Augustine, 
 renowned for his zeal for the faith, for holiness of life, and for miracles. 
 His feast is celebrated on the day following.
\switchcolumn*
\selectlanguage{latin}
Aquiléjæ pássio 
 sanctórum Felícis et Fortunáti fratrum, qui, in persecutióne Diocletiáni et 
 Maximiáni, equúleo suspénsi, atque, ardéntibus lampádibus circa eórum látera 
 appósitis et mox divína virtúte exstínctis, per ventrem fervénti óleo sunt 
 perfúsi; et ad últimum, cum in Christi confessióne persísterent, gloriósi 
 certáminis cursum, obtruncáti cápite, implevérunt.
\switchcolumn
\selectlanguage{english}
At Aquileia, the martyrdom of the 
 Saints Felix and Fortunatus, brothers. In the persecution of 
 Diocletian and Maximian, they were placed on the rack, and had flaming 
 torches held against their sides. These were extinguished by the power 
 of God, and boiling oil was poured over them. As they persevered in 
 confessing Christ, they were beheaded.
\switchcolumn*
\selectlanguage{latin}
Bremæ natális sancti 
 Rembérti, Hamburgénsis ac Breménsis Epíscopi.
\switchcolumn
\selectlanguage{english}
At Bremen, the birthday of St. 
 Rembert, bishop of Hamburg and Bremen.
\switchcolumn*
\selectlanguage{latin}
Tarvísii sancti Parísii, 
 civis Bononiénsis, Confessóris et Mónachi, ex Ordine Camaldulénsi.
\switchcolumn
\selectlanguage{english}
At Treviso, St. Parisius, a citizen 
 of Bologna, confessor and monk of the Camaldolese Order.
\switchcolumn*
\selectlanguage{latin}
Romæ Translátio sancti 
 Gregórii Nazianzéni, Epíscopi, Confessóris atque Ecclésiæ Doctóris; cujus 
 sacrum corpus, e Constantinópoli ántea delátum ad Urbem, atque in Ecclésia 
 sanctæ Dei Genitrícis ad Campum Mártium diu asservátum, Gregórius Décimus 
 tértius, Póntifex Máximus, in Capéllum, a se in Basílica sancti Petri 
 magnificentíssime exornátum, summa celebritáte tránstulit, ac postrídie 
 digno honóre sub altári cóndidit.
\switchcolumn
\selectlanguage{english}
At Rome, the translation of St. 
 Gregory Nazianzen, whose revered body was brought from Constantinople to 
 Rome, and kept for a long time in the Church of the Mother of God. It 
 was then transferred with great solemnity by Pope Gregory XIII to a chapel 
 of the basilica of St. Peter, magnificently decorated by His Holiness, and 
 the next day placed with due honour beneath the altar.
\switchcolumn*
\selectlanguage{latin}
\end{paracol}


% ---- martyrology/mart06/mart0612.htm
\needspace{10\baselineskip}
\begin{paracol}{2}
\selectlanguage{latin}
\begin{center}{\color{gregoriocolor} Prídie Idus Júnii. 
 Luna\dots\ }\end{center}
\switchcolumn
\selectlanguage{english}
\begin{center}{\color{gregoriocolor} The 12\textsuperscript{th} Day of 
 June. The\dots\ Day of the Moon.}\end{center}
\end{paracol}

\noindent\begin{tabularx}{\linewidth}{*{19}{>{\centering\arraybackslash}X}}
 \textcolor{gregoriocolor}{a} & \textcolor{gregoriocolor}{b} & \textcolor{gregoriocolor}{c} & \textcolor{gregoriocolor}{d} & \textcolor{gregoriocolor}{e} & \textcolor{gregoriocolor}{f} & \textcolor{gregoriocolor}{g} & \textcolor{gregoriocolor}{h} & \textcolor{gregoriocolor}{i} & \textcolor{gregoriocolor}{k} & \textcolor{gregoriocolor}{l} & \textcolor{gregoriocolor}{m} & \textcolor{gregoriocolor}{n} & \textcolor{gregoriocolor}{p} & \textcolor{gregoriocolor}{q} & \textcolor{gregoriocolor}{r} & \textcolor{gregoriocolor}{s} & \textcolor{gregoriocolor}{t} & \textcolor{gregoriocolor}{u} \\
 16 & 17 & 18 & 19 & 20 & 21 & 22 & 23 & 24 & 25 & 26 & 27 & 28 & 29 & 1 & 2 & 3 & 4 & 5 \\
\end{tabularx}
\vspace{0.5\baselineskip}
\noindent\begin{tabularx}{\linewidth}{*{12}{>{\centering\arraybackslash}X}}
 \textcolor{gregoriocolor}{A} & \textcolor{gregoriocolor}{B} & \textcolor{gregoriocolor}{C} & \textcolor{gregoriocolor}{D} & \textcolor{gregoriocolor}{E} & F & \textcolor{gregoriocolor}{F} & \textcolor{gregoriocolor}{G} & \textcolor{gregoriocolor}{H} & \textcolor{gregoriocolor}{M} & \textcolor{gregoriocolor}{N} & \textcolor{gregoriocolor}{P} \\
 6 & 7 & 8 & 9 & 10 & 11 & 10 & 11 & 12 & 13 & 14 & 15 \\
\end{tabularx}

\begin{paracol}{2}
\selectlanguage{latin}
\lettrine[lines=2]{S}{ancti} Joánnis a sancto Facúndo, ex Eremitárum sancti 
 Augustíni Ordine, Confessóris, qui migrávit in cælum prídie hujus diéi.
\switchcolumn
\selectlanguage{english}
\lettrine[lines=2]{S}{t.} John of St. Facundus, confessor of the Order of the 
 Hermits of St. Augustine, who died on the 11th of June.
\switchcolumn*
\selectlanguage{latin}
Romæ, via Aurélia, natális sanctórum Mártyrum Basílidis, 
 Cyríni, Náboris et Nazárii militum, qui, in persecutióne Diocletiáni et 
 Maximiáni, sub Aurélio Præfécto, ob Christiáni nóminis confessiónem, detrúsi 
 in cárcerem et scorpiónibus laceráti, tandem cápite truncáti sunt.
\switchcolumn
\selectlanguage{english}
At Rome, on the Aurelian Way, during the persecution of 
 Diocletian and Maximian, and under the prefect Aurelius, the birthday of the 
 holy martyrs Basilides, Cyrinus, Nabor, and Nazarius, all soldiers who were 
 cast into prison for the confession of the Christian name, scourged with 
 knotted whips, and finally beheaded.
\switchcolumn*
\selectlanguage{latin}
Nicǽæ, in Bithynia, sanctæ Antonínæ Mártyris, quæ, in 
 eádem persecutióne, a Priscilliáno Præside jussa est fústibus cædi, suspéndi 
 in equúleo, latéribus laniári et flammis incéndi, ac demum gládio necári.
\switchcolumn
\selectlanguage{english}
At Nicaea in Bithynia, St. Antonina, martyr. She 
 was scourged by order of the governor Priscillian during the same 
 persecution, then racked, lacerated, exposed to the fire, and finally put to 
 the sword.
\switchcolumn*
\selectlanguage{latin}
Romæ, in Basílica Vaticána, sancti Leónis Papæ Tértii, 
 cui erútos ab ímpiis óculos et præcísam linguam Deus mirabíliter restítuit.
\switchcolumn
\selectlanguage{english}
At Rome, in the Vatican basilica, Pope St. Leo III, to 
 whom God miraculously restored his eyes and his tongue after they had been 
 torn out by impious men.
\switchcolumn*
\selectlanguage{latin}
In Thrácia sancti Olympii Epíscopi, qui, ab Ariánis sede 
 pulsus, Conféssor occúbuit.
\switchcolumn
\selectlanguage{english}
In Thrace, St. Olympius, a bishop, who was driven out of 
 his diocese by the Arians, and died a confessor.
\switchcolumn*
\selectlanguage{latin}
In Cilícia sancti Amphiónis Epíscopi, qui, témpore 
 Galérii Maximiáni, Conféssor fuit egrégius.
\switchcolumn
\selectlanguage{english}
In Cilicia, Bishop St. Amphion, a celebrated confessor of 
 the time of Galerius Maximian.
\switchcolumn*
\selectlanguage{latin}
In Ægypto sancti Onúphrii Anachorétæ, qui in vasta erémo 
 sexagínta annis vitam religióse perégit, et, magnis virtútibus ac méritis 
 clarus, migrávit in cælum. Ipsíus vero insígnia gesta Paphnútius Abbas 
 conscrípsit.
\switchcolumn
\selectlanguage{english}
In Egypt, St. Onuphrius, an anchoret, who for sixty years 
 led a religious life in the desert, and renowned for great virtues and 
 merits departed for heaven. His admirable deeds have been recorded by 
 Abbot Paphnutius.
\switchcolumn*
\selectlanguage{latin}
\end{paracol}


% ---- martyrology/mart06/mart0613.htm
\needspace{10\baselineskip}
\begin{paracol}{2}
\selectlanguage{latin}
\begin{center}{\color{gregoriocolor} Idibus Júnii. 
 Luna\dots\ }\end{center}
\switchcolumn
\selectlanguage{english}
\begin{center}{\color{gregoriocolor} The 13\textsuperscript{th} Day of 
 June. The\dots\ Day of the Moon.}\end{center}
\end{paracol}

\noindent\begin{tabularx}{\linewidth}{*{19}{>{\centering\arraybackslash}X}}
 \textcolor{gregoriocolor}{a} & \textcolor{gregoriocolor}{b} & \textcolor{gregoriocolor}{c} & \textcolor{gregoriocolor}{d} & \textcolor{gregoriocolor}{e} & \textcolor{gregoriocolor}{f} & \textcolor{gregoriocolor}{g} & \textcolor{gregoriocolor}{h} & \textcolor{gregoriocolor}{i} & \textcolor{gregoriocolor}{k} & \textcolor{gregoriocolor}{l} & \textcolor{gregoriocolor}{m} & \textcolor{gregoriocolor}{n} & \textcolor{gregoriocolor}{p} & \textcolor{gregoriocolor}{q} & \textcolor{gregoriocolor}{r} & \textcolor{gregoriocolor}{s} & \textcolor{gregoriocolor}{t} & \textcolor{gregoriocolor}{u} \\
 17 & 18 & 19 & 20 & 21 & 22 & 23 & 24 & 25 & 26 & 27 & 28 & 29 & 1 & 2 & 3 & 4 & 5 & 6 \\
\end{tabularx}
\vspace{0.5\baselineskip}
\noindent\begin{tabularx}{\linewidth}{*{12}{>{\centering\arraybackslash}X}}
 \textcolor{gregoriocolor}{A} & \textcolor{gregoriocolor}{B} & \textcolor{gregoriocolor}{C} & \textcolor{gregoriocolor}{D} & \textcolor{gregoriocolor}{E} & F & \textcolor{gregoriocolor}{F} & \textcolor{gregoriocolor}{G} & \textcolor{gregoriocolor}{H} & \textcolor{gregoriocolor}{M} & \textcolor{gregoriocolor}{N} & \textcolor{gregoriocolor}{P} \\
 7 & 8 & 9 & 10 & 11 & 12 & 11 & 12 & 13 & 14 & 15 & 16 \\
\end{tabularx}

\begin{paracol}{2}
\selectlanguage{latin}
\lettrine[lines=2]{P}{atávii} sancti Antónii 
 Lusitáni, Sacerdótis ex Ordine Minórum et Confessóris, atque Ecclésiæ 
 Doctóris, vita et miráculis, ac prædicatióne illústris, quem, uno post illíus óbitum anno nondum expléto, Gregórius Papa Nonus in Sanctórum cánonem 
 rétulit.
\switchcolumn
\selectlanguage{english}
\lettrine[lines=2]{A}{t} Padua, St. Anthony, a native of 
 Portugal, priest of the Order of Friars Minor and confessor, illustrious for 
 the sanctity of his life, his miracles, and his preaching. Pope 
 Gregory IX placed him on the canon of the saints within a year after his 
 death.
\switchcolumn*
\selectlanguage{latin}
Romæ, via Ardeatína, 
 natális sanctæ Felículæ, Vírginis et Mártyris; quæ, nec Flacco núbere neque 
 idólis immoláre volens, trádita est cuídam Júdici, qui eam in confessióne 
 Christi perseverántem, post tenebricósam custódiam et famis inédiam, támdiu 
 fecit in equúleo torquéri, donec illa emítteret spíritum, et sic demum 
 depóni et in cloácam præcipitári. Ipsíus vero corpus, inde extráctum, 
 sanctus Nicomédes Présbyter eádem via sepelívit.
\switchcolumn
\selectlanguage{english}
At Rome, on the Ardeatine Way, the 
 birthday of St. Felicula, virgin and martyr, who was delivered to the judge 
 for refusing to marry Flaccus and to sacrifice to idols. As she 
 persevered in the confession of Christ, he confined her in a dark dungeon 
 without food, and afterwards caused her to be stretched on the rack until 
 she expired. She was then thrown into a sewer, but St. Nicomedes the 
 Priest recovered her body and buried it on this road.
\switchcolumn*
\selectlanguage{latin}
In Pelígnis sancti 
 Peregríni, Epíscopi et Mártyris, qui, pro fide cathólica, a Longobárdis in 
 Atérnum flumen demérsus est.
\switchcolumn
\selectlanguage{english}
In Abruzzi, St. Peregrinus, bishop 
 and martyr. For the Catholic faith he was thrown into the river Aterno 
 by the Lombards.
\switchcolumn*
\selectlanguage{latin}
Córdubæ, in Hispánia, 
 sancti Fándilæ, Presbyteri et Mónachi; qui, in persecutióne Arábica, 
 amputáto cápite, pro Christi fide martyrium súbiit.
\switchcolumn
\selectlanguage{english}
At Cordova in Spain, in the 
 persecution of the Arabs, St. Fandila, a priest and monk, who underwent 
 martyrdom by beheading for the faith of Christ.
\switchcolumn*
\selectlanguage{latin}
In Africa sanctórum 
 Mártyrum Fortunáti et Luciáni.
\switchcolumn
\selectlanguage{english}
In Africa, the holy martyrs 
 Fortunatus and Lucian.
\switchcolumn*
\selectlanguage{latin}
Bybli, in Phœnícia, 
 sanctæ Aquilínæ, Vírginis et Mártyris, quæ, annos duódecim nata, sub 
 Diocletiáno Imperatóre et Volusiáno Júdice, ob fídei confessiónem cólaphis 
 et verbéribus cæsa, et súbulis candéntibus perforáta, demum, percússa gládio, 
 virginitátem martyrio consecrávit.
\switchcolumn
\selectlanguage{english}
At Byblos in Phoenicia, St. Aquilina, 
 virgin and martyr, at the age of twelve years, under Emperor Diocletian and 
 the judge Volusian. For the confession of the faith, she was beaten, 
 scourged, pierced with heated stakes, and finally being struck with a sword, 
 consecrated her virginity by martyrdom.
\switchcolumn*
\selectlanguage{latin}
In Cypro sancti 
 Triphylii Epíscopi.
\switchcolumn
\selectlanguage{english}
In Cyprus, St. Triphyllius, bishop.
\switchcolumn*
\selectlanguage{latin}
\end{paracol}


% ---- martyrology/mart06/mart0614.htm
\needspace{10\baselineskip}
\begin{paracol}{2}
\selectlanguage{latin}
\begin{center}{\color{gregoriocolor} Décimo octávo Kaléndas Júlii. 
 Luna\dots\ }\end{center}
\switchcolumn
\selectlanguage{english}
\begin{center}{\color{gregoriocolor} The 14\textsuperscript{th} Day of 
 June. The\dots\ Day of the Moon.}\end{center}
\end{paracol}

\noindent\begin{tabularx}{\linewidth}{*{19}{>{\centering\arraybackslash}X}}
 \textcolor{gregoriocolor}{a} & \textcolor{gregoriocolor}{b} & \textcolor{gregoriocolor}{c} & \textcolor{gregoriocolor}{d} & \textcolor{gregoriocolor}{e} & \textcolor{gregoriocolor}{f} & \textcolor{gregoriocolor}{g} & \textcolor{gregoriocolor}{h} & \textcolor{gregoriocolor}{i} & \textcolor{gregoriocolor}{k} & \textcolor{gregoriocolor}{l} & \textcolor{gregoriocolor}{m} & \textcolor{gregoriocolor}{n} & \textcolor{gregoriocolor}{p} & \textcolor{gregoriocolor}{q} & \textcolor{gregoriocolor}{r} & \textcolor{gregoriocolor}{s} & \textcolor{gregoriocolor}{t} & \textcolor{gregoriocolor}{u} \\
 18 & 19 & 20 & 21 & 22 & 23 & 24 & 25 & 26 & 27 & 28 & 29 & 1 & 2 & 3 & 4 & 5 & 6 & 7 \\
\end{tabularx}
\vspace{0.5\baselineskip}
\noindent\begin{tabularx}{\linewidth}{*{12}{>{\centering\arraybackslash}X}}
 \textcolor{gregoriocolor}{A} & \textcolor{gregoriocolor}{B} & \textcolor{gregoriocolor}{C} & \textcolor{gregoriocolor}{D} & \textcolor{gregoriocolor}{E} & F & \textcolor{gregoriocolor}{F} & \textcolor{gregoriocolor}{G} & \textcolor{gregoriocolor}{H} & \textcolor{gregoriocolor}{M} & \textcolor{gregoriocolor}{N} & \textcolor{gregoriocolor}{P} \\
 8 & 9 & 10 & 11 & 12 & 13 & 12 & 13 & 14 & 15 & 16 & 17 \\
\end{tabularx}

\begin{paracol}{2}
\selectlanguage{latin}
\lettrine[lines=2]{S}{ancti} Basilíi, 
 cognoménto Magni, Confessóris et Ecclésiæ Doctóris, qui Kaléndis Januárii 
 obdormívit in Dómino, sed hac die potíssimum cólitur, qua Cæsariénsis in 
 Cappadócia Epíscopus ordinátus est.
\switchcolumn
\selectlanguage{english}
\lettrine[lines=2]{S}{t.} Basil, surnamed the Great, 
 confessor and doctor of the Church. He died on the 1st of January, but 
 his feast is celebrated today, for it was on this day that he was 
 consecrated bishop of Caesarea in Cappadocia.
\switchcolumn*
\selectlanguage{latin}
Samaríæ, in Palæstína, 
 sancti Eliséi Prophétæ, cujus sepúlcrum, ubi et Prophéta quiéscit Abdías, a 
 dæmónibus perhorrésci sanctus Hierónymus scribit.
\switchcolumn
\selectlanguage{english}
At Samaria in Palestine, the holy 
 prophet Eliseus, whose grave, says St. Jerome, makes the demons tremble. 
 With him also rests the prophet Abdias.
\switchcolumn*
\selectlanguage{latin}
Syracúsis, in Sicília, 
 sancti Marciáni Epíscopi, qui, a beáto Petro Apóstolo ordinátus Epíscopus, 
 ibídem, post Evangélii prædicatiónem, a Judæis occísus est.
\switchcolumn
\selectlanguage{english}
At Syracuse in Sicily, St. Marcian, 
 bishop, who was made bishop by blessed Peter, and killed by the Jews after 
 he had preached the Gospel.
\switchcolumn*
\selectlanguage{latin}
Córdubæ, in Hispánia, 
 sanctórum Mártyrum Anastásii Presbyteri, Felícis Mónachi, et Dignæ Vírginis.
\switchcolumn
\selectlanguage{english}
At Cordova in Spain, the holy 
 martyrs Anastasius, a priest, Felix, a monk, and Digna, virgin.
\switchcolumn*
\selectlanguage{latin}
Apud Suessiónes, in 
 Gálliis, sanctórum mártyrum Valérii et Rufíni, qui, in persecutióne 
 Diocletiáni, a Præside Rictiováro, post multa torménta, jussi sunt decollári.
\switchcolumn
\selectlanguage{english}
At Soissons in France, the holy 
 martyrs Valerius and Rufinus, who, after enduring many torments, were 
 condemned to be beheaded by the governor Rictiovarus, in the persecution of 
 Diocletian.
\switchcolumn*
\selectlanguage{latin}
Constantinópoli sancti 
 Methódii Epíscopi.
\switchcolumn
\selectlanguage{english}
At Constantinople, St. Methodius, 
 bishop.
\switchcolumn*
\selectlanguage{latin}
Viénnæ, in Gállia, 
 sancti Æthérii Epíscopi.
\switchcolumn
\selectlanguage{english}
At Vienne, St. Aetherius, bishop.
\switchcolumn*
\selectlanguage{latin}
Apud Rhuténos, in 
 Gállia, sancti Quinctiáni Epíscopi.
\switchcolumn
\selectlanguage{english}
At Rodez in France, St. Quinctian, 
 bishop.
\switchcolumn*
\selectlanguage{latin}
\end{paracol}


% ---- martyrology/mart06/mart0615.htm
\needspace{10\baselineskip}
\begin{paracol}{2}
\selectlanguage{latin}
\begin{center}{\color{gregoriocolor} Décimo séptimo Kaléndas Júlii. 
 Luna\dots\ }\end{center}
\switchcolumn
\selectlanguage{english}
\begin{center}{\color{gregoriocolor} The 15\textsuperscript{th} Day of 
 June. The\dots\ Day of the Moon.}\end{center}
\end{paracol}

\noindent\begin{tabularx}{\linewidth}{*{19}{>{\centering\arraybackslash}X}}
 \textcolor{gregoriocolor}{a} & \textcolor{gregoriocolor}{b} & \textcolor{gregoriocolor}{c} & \textcolor{gregoriocolor}{d} & \textcolor{gregoriocolor}{e} & \textcolor{gregoriocolor}{f} & \textcolor{gregoriocolor}{g} & \textcolor{gregoriocolor}{h} & \textcolor{gregoriocolor}{i} & \textcolor{gregoriocolor}{k} & \textcolor{gregoriocolor}{l} & \textcolor{gregoriocolor}{m} & \textcolor{gregoriocolor}{n} & \textcolor{gregoriocolor}{p} & \textcolor{gregoriocolor}{q} & \textcolor{gregoriocolor}{r} & \textcolor{gregoriocolor}{s} & \textcolor{gregoriocolor}{t} & \textcolor{gregoriocolor}{u} \\
 19 & 20 & 21 & 22 & 23 & 24 & 25 & 26 & 27 & 28 & 29 & 1 & 2 & 3 & 4 & 5 & 6 & 7 & 8 \\
\end{tabularx}
\vspace{0.5\baselineskip}
\noindent\begin{tabularx}{\linewidth}{*{12}{>{\centering\arraybackslash}X}}
 \textcolor{gregoriocolor}{A} & \textcolor{gregoriocolor}{B} & \textcolor{gregoriocolor}{C} & \textcolor{gregoriocolor}{D} & \textcolor{gregoriocolor}{E} & F & \textcolor{gregoriocolor}{F} & \textcolor{gregoriocolor}{G} & \textcolor{gregoriocolor}{H} & \textcolor{gregoriocolor}{M} & \textcolor{gregoriocolor}{N} & \textcolor{gregoriocolor}{P} \\
 9 & 10 & 11 & 12 & 13 & 14 & 13 & 14 & 15 & 16 & 17 & 18 \\
\end{tabularx}

\begin{paracol}{2}
\selectlanguage{latin}
\lettrine[lines=2]{A}{pud} Silárum flumen, in 
 Lucánia, natális sanctórum Mártyrum Viti, Modésti atque Crescéntiæ, qui, sub 
 Diocletiáno Imperatóre, e Sicília illuc deláti, ibídem, post ollam fervéntis 
 plumbi, post béstias et catástas divína virtúte superátas, cursum gloriósi 
 certáminis peregérunt.
\switchcolumn
\selectlanguage{english}
\lettrine[lines=2]{N}{ear} the river Silaro in Lucania, 
 the birthday of the holy martyrs Vitus, Modestus, and Crescentia, who were 
 brought there from Sicily in the reign of the emperor Diocletian. They 
 were plunged into a vessel of molten lead, exposed to the beasts, and 
 stretched on the rack, but after having survived these torments through the 
 power of God, they came to the end of their glorious trials.
\switchcolumn*
\selectlanguage{latin}
Doróstori, in Mysia 
 inferióre, sancti Hesychii mílitis, qui, cum beáto Júlio comprehénsus, post 
 eum, sub Máximo Præside, martyrio coronátus est.
\switchcolumn
\selectlanguage{english}
At Silistria in Rumania, St. 
 Hesychius, a soldier, who was arrested with blessed Julius, and under the 
 governor Máximus followed him to the crown of martyrdom.
\switchcolumn*
\selectlanguage{latin}
Zephyrii, in Cilícia, 
 sancti Dulæ Mártyris, qui sub Præside Máximo, ob Christi nomen, virgis cæsus, 
 in cratícula pósitus, fervénti óleo incénsus aliáque passus, martyrii palmam 
 victor accépit.
\switchcolumn
\selectlanguage{english}
At Zephirium in Cilicia, St. Dulas, 
 martyr under the governor Maximus. For the name of Christ, he was 
 scourged, laid on the gridiron, scalded with boiling oil, and after enduring 
 other trials, received for his victory the palm of martyrdom.
\switchcolumn*
\selectlanguage{latin}
Córdubæ, in Hispánia, 
 sanctæ Beníldis Mártyris.
\switchcolumn
\selectlanguage{english}
At Cordova in Spain, St. Benildes, 
 martyr.
\switchcolumn*
\selectlanguage{latin}
Sibápoli, in 
 Mesopotámia, sanctárum Vírginum et Mártyrum Lybes et Leónidis sorórum, et 
 Eutrópiæ, puéllæ annórum duódecim; quæ per divérsa torménta ad corónam 
 martyrii pervenérunt.
\switchcolumn
\selectlanguage{english}
At Palmyra in Sicily, the holy 
 martyrs Libya and Leonides, sisters, and Eutropia, a girl of twelve years, 
 who won the crown of martyrdom by various torments.
\switchcolumn*
\selectlanguage{latin}
Apud Valencénas, in 
 Gállia, deposítio sancti Landelíni Abbátis.
\switchcolumn
\selectlanguage{english}
At Valenciennes in France, the 
 death of St. Landelin, abbot.
\switchcolumn*
\selectlanguage{latin}
Arvérnis, in Gállia, 
 sancti Abrahæ Confessóris, sanctitáte ac miráculis illústris.
\switchcolumn
\selectlanguage{english}
In Auvergne in France, St. Abraham, 
 confessor, illustrious by his holiness and miracles.
\switchcolumn*
\selectlanguage{latin}
Pibráci, in diœcési 
 Tolosáni, sanctæ Germánæ Cousin, Vírginis, quæ custodiéndis grégibus addícta 
 permánsit, et, cum húmilis ac pauper vixísset, et, multis ærúmnis 
 patientíssime tolerátis, migrásset ad Sponsum, plúrimis post óbitum 
 miráculis cláruit; atque a Pio Nono, Pontífice Máximo, in sanctárum Vírginum 
 número adscrípta fuit.
\switchcolumn
\selectlanguage{english}
At Pibrac in the diocese of 
 Toulouse, St. Germaine Cousin, virgin. After a life of poverty, 
 humility, and patient suffering amidst many trials as shepherdess of her 
 flocks, she went to her heavenly spouse, and became renowned for numerous 
 miracles after her death. Pope Pius IX placed her in the number of 
 holy virgins.
\switchcolumn*
\selectlanguage{latin}
\end{paracol}


% ---- martyrology/mart06/mart0616.htm
\needspace{10\baselineskip}
\begin{paracol}{2}
\selectlanguage{latin}
\begin{center}{\color{gregoriocolor} Sextodécimo Kaléndas Júlii. 
 Luna\dots\ }\end{center}
\switchcolumn
\selectlanguage{english}
\begin{center}{\color{gregoriocolor} The 16\textsuperscript{th} Day of 
 June. The\dots\ Day of the Moon.}\end{center}
\end{paracol}

\noindent\begin{tabularx}{\linewidth}{*{19}{>{\centering\arraybackslash}X}}
 \textcolor{gregoriocolor}{a} & \textcolor{gregoriocolor}{b} & \textcolor{gregoriocolor}{c} & \textcolor{gregoriocolor}{d} & \textcolor{gregoriocolor}{e} & \textcolor{gregoriocolor}{f} & \textcolor{gregoriocolor}{g} & \textcolor{gregoriocolor}{h} & \textcolor{gregoriocolor}{i} & \textcolor{gregoriocolor}{k} & \textcolor{gregoriocolor}{l} & \textcolor{gregoriocolor}{m} & \textcolor{gregoriocolor}{n} & \textcolor{gregoriocolor}{p} & \textcolor{gregoriocolor}{q} & \textcolor{gregoriocolor}{r} & \textcolor{gregoriocolor}{s} & \textcolor{gregoriocolor}{t} & \textcolor{gregoriocolor}{u} \\
 20 & 21 & 22 & 23 & 24 & 25 & 26 & 27 & 28 & 29 & 1 & 2 & 3 & 4 & 5 & 6 & 7 & 8 & 9 \\
\end{tabularx}
\vspace{0.5\baselineskip}
\noindent\begin{tabularx}{\linewidth}{*{12}{>{\centering\arraybackslash}X}}
 \textcolor{gregoriocolor}{A} & \textcolor{gregoriocolor}{B} & \textcolor{gregoriocolor}{C} & \textcolor{gregoriocolor}{D} & \textcolor{gregoriocolor}{E} & F & \textcolor{gregoriocolor}{F} & \textcolor{gregoriocolor}{G} & \textcolor{gregoriocolor}{H} & \textcolor{gregoriocolor}{M} & \textcolor{gregoriocolor}{N} & \textcolor{gregoriocolor}{P} \\
 10 & 11 & 12 & 13 & 14 & 15 & 14 & 15 & 16 & 17 & 18 & 19 \\
\end{tabularx}

\begin{paracol}{2}
\selectlanguage{latin}
\lettrine[lines=2]{M}{ogúntiæ} pássio 
 sanctórum Auræi Epíscopi, et Justínæ soróris, ac ceterórum Mártyrum; qui, 
 synáxim in Ecclésia agéntes, ab Hunnis, Germániam diripiéntibus, trucidáti 
 sunt.
\switchcolumn
\selectlanguage{english}
\lettrine[lines=2]{A}{t} Mainz, the passion of the Saints 
 Aureus and Justina, his sister, and other martyrs who were massacred by the 
 Huns, at that time devastating Germany, while they were in church at Mass.
\switchcolumn*
\selectlanguage{latin}
Vesontióne, in Gálliis, 
 sanctórum Mártyrum Ferreóli Presbyteri, et Ferruntiónis Diáconi, qui, a 
 beáto Irenæo Epíscopo missi ad prædicándum verbum Dei, póstea, sub Cláudio 
 Júdice, divérsis pœnis excruciáti, gládio feriúntur.
\switchcolumn
\selectlanguage{english}
At Besançon in France, the holy 
 martyrs Ferreol, a priest, and Ferruntion, a deacon, who were sent by the 
 blessed bishop Irenæus to preach the word of God, and after being exposed 
 to various torments under Judge Claudius, were put to the sword.
\switchcolumn*
\selectlanguage{latin}
Tarsi, in Cilícia, 
 sanctórum Mártyrum Quírici et Julíttæ matris, sub Diocletiáno Imperatóre. 
 Ex eis Quíricus, triénnis puéllus, cum matrem, quæ ante Alexándrum Præsidem 
 nervis diríssime cædebátur, implacábili luctu lugéret, ad gradus tribunális 
 illísus intériit; Julítta vero, post dira vérbera et grávia torménta, 
 martyrii sui cursum obtruncatióne cápitis implévit.
\switchcolumn
\selectlanguage{english}
At Tarsus in Cilicia, in the reign 
 of Emperor Diocletian, the holy martyrs Cyricus and Julitta, his mother. 
 Cyricus, a child of three years, seeing his mother cruelly scourged with 
 whips in the presence of the governor Alexander, and crying bitterly, was 
 killed by being dashed against the steps of the tribunal. Julitta, 
 after being subjected to severe lashings and grievous torments, closed the 
 course of her martyrdom by beheading.
\switchcolumn*
\selectlanguage{latin}
Amathúnte, in Cypro, 
 sancti Tychónis Epíscopi, témpore Theodósii junióris.
\switchcolumn
\selectlanguage{english}
At Amathus in Cyprus, St. Tychon, a 
 bishop in the time of Theodosius the Younger.
\switchcolumn*
\selectlanguage{latin}
Lugdúni, in Gállia, 
 deposítio beáti Aureliáni, Epíscopi Arelaténsis.
\switchcolumn
\selectlanguage{english}
At Lyons, the death of blessed 
 Aurelian, bishop of Arles.
\switchcolumn*
\selectlanguage{latin}
Nannéte, in Británnia minóre, sancti Similiáni, Epíscopi et Confessóris.
\switchcolumn
\selectlanguage{english}
At Nantes in Brittany, St. Similian, 
 bishop and confessor.
\switchcolumn*
\selectlanguage{latin}
Misnæ, in Germánia, 
 sancti Bennónis Epíscopi.
\switchcolumn
\selectlanguage{english}
At Meissen in Germany, St. Benno, 
 bishop.
\switchcolumn*
\selectlanguage{latin}
In monastério 
 Aquiriénsi, in Brabántia, sanctæ Lutgárdis Vírginis.
\switchcolumn
\selectlanguage{english}
In the monastery of Aywieres in 
 Brabant, St. Lutgard, virgin.
\switchcolumn*
\selectlanguage{latin}
\end{paracol}


% ---- martyrology/mart06/mart0617.htm
\needspace{10\baselineskip}
\begin{paracol}{2}
\selectlanguage{latin}
\begin{center}{\color{gregoriocolor} Quintodécimo Kaléndas Júlii. 
 Luna\dots\ }\end{center}
\switchcolumn
\selectlanguage{english}
\begin{center}{\color{gregoriocolor} The 17\textsuperscript{th} Day of 
 June. The\dots\ Day of the Moon.}\end{center}
\end{paracol}

\noindent\begin{tabularx}{\linewidth}{*{19}{>{\centering\arraybackslash}X}}
 \textcolor{gregoriocolor}{a} & \textcolor{gregoriocolor}{b} & \textcolor{gregoriocolor}{c} & \textcolor{gregoriocolor}{d} & \textcolor{gregoriocolor}{e} & \textcolor{gregoriocolor}{f} & \textcolor{gregoriocolor}{g} & \textcolor{gregoriocolor}{h} & \textcolor{gregoriocolor}{i} & \textcolor{gregoriocolor}{k} & \textcolor{gregoriocolor}{l} & \textcolor{gregoriocolor}{m} & \textcolor{gregoriocolor}{n} & \textcolor{gregoriocolor}{p} & \textcolor{gregoriocolor}{q} & \textcolor{gregoriocolor}{r} & \textcolor{gregoriocolor}{s} & \textcolor{gregoriocolor}{t} & \textcolor{gregoriocolor}{u} \\
 21 & 22 & 23 & 24 & 25 & 26 & 27 & 28 & 29 & 1 & 2 & 3 & 4 & 5 & 6 & 7 & 8 & 9 & 10 \\
\end{tabularx}
\vspace{0.5\baselineskip}
\noindent\begin{tabularx}{\linewidth}{*{12}{>{\centering\arraybackslash}X}}
 \textcolor{gregoriocolor}{A} & \textcolor{gregoriocolor}{B} & \textcolor{gregoriocolor}{C} & \textcolor{gregoriocolor}{D} & \textcolor{gregoriocolor}{E} & F & \textcolor{gregoriocolor}{F} & \textcolor{gregoriocolor}{G} & \textcolor{gregoriocolor}{H} & \textcolor{gregoriocolor}{M} & \textcolor{gregoriocolor}{N} & \textcolor{gregoriocolor}{P} \\
 11 & 12 & 13 & 14 & 15 & 16 & 15 & 16 & 17 & 18 & 19 & 20 \\
\end{tabularx}

\begin{paracol}{2}
\selectlanguage{latin}
\lettrine[lines=2]{R}{omæ} natális sanctórum 
 ducentórum et sexagínta duórum Mártyrum, qui, in persecutióne Diocletiáni, 
 pro fide Christi sunt necáti, ac pósiti via Salária véteri, ad clivum 
 Cucúmeris.
\switchcolumn
\selectlanguage{english}
\lettrine[lines=2]{A}{t} Rome, during the persecution of 
 Diocletian, the birthday of two hundred and sixty-two martyrs, who were put 
 to death for the faith of Christ, and buried on the old Salarian Way, at the 
 foot of Cucumer Hill.
\switchcolumn*
\selectlanguage{latin}
Vesontióne, in Gálliis, 
 sancti Antídii, Epíscopi et Mártyris, qui, ob Christi fidem, a Wándalis 
 occísus fuit.
\switchcolumn
\selectlanguage{english}
At Besançon in France, St. Antidius, 
 bishop and martyr, who was slain by the Vandals for the faith of Christ.
\switchcolumn*
\selectlanguage{latin}
Apollóniæ, in 
 Macedónia, sanctórum Mártyrum Atheniénsium Isáuri Diáconi, Innocéntii, 
 Felícis, Jeremíæ et Peregríni, qui, a Tripóntio Tribúno várie torti, cápite 
 obtruncáti sunt.
\switchcolumn
\selectlanguage{english}
At Apollonia in Macedonia, the holy 
 martyrs Isaurus, a deacon, Innocent, Felix, Jeremias, and Peregrínus, all of 
 them Athenians who were tortured in various ways by the tribune Tripontius, 
 and beheaded.
\switchcolumn*
\selectlanguage{latin}
Tarracínæ, in Campánia, 
 sancti Montáni mílitis, qui sub Hadriáno Imperatóre et Leóntio Consulári, 
 post multa torménta, martyrii corónam accépit.
\switchcolumn
\selectlanguage{english}
At Terracina in Campania, St. 
 Montanus, a soldier, who received the crown of martyrdom after suffering 
 many torments, in the time of Emperor Hadrian and the governor Leontius.
\switchcolumn*
\selectlanguage{latin}
Apud Venáfrum, in 
 Campánia, sanctórum Mártyrum Nicándri et Marciáni, qui, in persecutióne 
 Maximiáni, cápite cæsi sunt.
\switchcolumn
\selectlanguage{english}
At Venafro in Campania, the holy 
 martyrs Nicander and Marcian, who were beheaded in the persecution of 
 Maximian.
\switchcolumn*
\selectlanguage{latin}
Chalcédone sanctórum 
 Mártyrum Manuélis, Sabélis et Ismaélis, qui, pacis causa apud Juliánum 
 Apóstatam pro Persárum Rege legatióne fungéntes, ab ipso Imperatóre, cum 
 idóla venerári compelleréntur idque constánti ánimo recusárent, gládio 
 feríri jubéntur.
\switchcolumn
\selectlanguage{english}
At Chalcedon, the holy martyrs 
 Manuel, Sabel, and Ismael, whom the king of Persia sent as ambassadors to 
 Julian the Apostate to obtain peace. Having firmly refused to worship 
 idols when commanded by the emperor, they were put to the sword.
\switchcolumn*
\selectlanguage{latin}
Amériæ, in Umbria, 
 sancti Himérii Epíscopi, cujus corpus Cremónam, in Insúbria, translátum est.
\switchcolumn
\selectlanguage{english}
At Amelia in Umbria, Bishop St. 
 Himerius, whose body was translated to Cremona.
\switchcolumn*
\selectlanguage{latin}
In pago Bituricénsi, 
 sancti Gundúlphi Epíscopi.
\switchcolumn
\selectlanguage{english}
In the territory of Bourges, St. 
 Gundulphus, bishop.
\switchcolumn*
\selectlanguage{latin}
Aureliánis, in Gállia, 
 sancti Avíti, Presbyteri et Confessóris.
\switchcolumn
\selectlanguage{english}
At Orleans in France, St. Avitus, 
 priest and confessor.
\switchcolumn*
\selectlanguage{latin}
In Phrygia sancti 
 Hypátii Confessóris.
\switchcolumn
\selectlanguage{english}
In Phrygia, St. Hypatius, confessor.
\switchcolumn*
\selectlanguage{latin}
Item sancti Bessariónis 
 Anachorétæ.
\switchcolumn
\selectlanguage{english}
Also, St. Bessarion, anchoret.
\switchcolumn*
\selectlanguage{latin}
Pisis, in Túscia, 
 sancti Rainérii Confessóris.
\switchcolumn
\selectlanguage{english}
At Pisa in Tuscany, St. Rainerius, 
 confessor.
\switchcolumn*
\selectlanguage{latin}
\end{paracol}


% ---- martyrology/mart06/mart0618.htm
\needspace{10\baselineskip}
\begin{paracol}{2}
\selectlanguage{latin}
\begin{center}{\color{gregoriocolor} Quartodécimo Kaléndas Júlii. 
 Luna\dots\ }\end{center}
\switchcolumn
\selectlanguage{english}
\begin{center}{\color{gregoriocolor} The 18\textsuperscript{th} Day of 
 June. The\dots\ Day of the Moon.}\end{center}
\end{paracol}

\noindent\begin{tabularx}{\linewidth}{*{19}{>{\centering\arraybackslash}X}}
 \textcolor{gregoriocolor}{a} & \textcolor{gregoriocolor}{b} & \textcolor{gregoriocolor}{c} & \textcolor{gregoriocolor}{d} & \textcolor{gregoriocolor}{e} & \textcolor{gregoriocolor}{f} & \textcolor{gregoriocolor}{g} & \textcolor{gregoriocolor}{h} & \textcolor{gregoriocolor}{i} & \textcolor{gregoriocolor}{k} & \textcolor{gregoriocolor}{l} & \textcolor{gregoriocolor}{m} & \textcolor{gregoriocolor}{n} & \textcolor{gregoriocolor}{p} & \textcolor{gregoriocolor}{q} & \textcolor{gregoriocolor}{r} & \textcolor{gregoriocolor}{s} & \textcolor{gregoriocolor}{t} & \textcolor{gregoriocolor}{u} \\
 22 & 23 & 24 & 25 & 26 & 27 & 28 & 29 & 1 & 2 & 3 & 4 & 5 & 6 & 7 & 8 & 9 & 10 & 11 \\
\end{tabularx}
\vspace{0.5\baselineskip}
\noindent\begin{tabularx}{\linewidth}{*{12}{>{\centering\arraybackslash}X}}
 \textcolor{gregoriocolor}{A} & \textcolor{gregoriocolor}{B} & \textcolor{gregoriocolor}{C} & \textcolor{gregoriocolor}{D} & \textcolor{gregoriocolor}{E} & F & \textcolor{gregoriocolor}{F} & \textcolor{gregoriocolor}{G} & \textcolor{gregoriocolor}{H} & \textcolor{gregoriocolor}{M} & \textcolor{gregoriocolor}{N} & \textcolor{gregoriocolor}{P} \\
 12 & 13 & 14 & 15 & 16 & 17 & 16 & 17 & 18 & 19 & 20 & 21 \\
\end{tabularx}

\begin{paracol}{2}
\selectlanguage{latin}
\lettrine[lines=2]{E}{déssæ,} in Mesopotámia, sancti Ephræm, Diáconi Edesséni et Confessóris, qui, post multos labóres pro 
 Christi fide suscéptos, doctrína et sanctitáte conspícuus, sub Valénte 
 Imperatóre quiévit in Dómino, et a Benedícto Papa Décimo quinto Doctor 
 Ecclésiæ universális est declarátus.
\switchcolumn
\selectlanguage{english}
\lettrine[lines=2]{A}{t} Edessa in Mesopotamia, St. 
 Ephraem, deacon of the church of Edessa in the time of Emperor Valens and 
 confessor. After suffering many trials for the faith of Christ and 
 gaining great renown for holiness and learning, he went to rest in the Lord. 
 He was declared a doctor of the Universal Church by Pope Benedict XV
\switchcolumn*
\selectlanguage{latin}
Romæ, via Ardeatína, 
 natális sanctórum Mártyrum Marci et Marcelliáni fratrum, qui, in 
 persecutióne Diocletiáni, a Júdice Fabiáno tenti et ad stípitem alligáti, in 
 pédibus clavos auctos accepérunt; cumque non cessárent laudáre Christum, 
 lánceis per látera transfíxi sunt, et cum glória martyrii ad sidérea regna 
 migrárunt.
\switchcolumn
\selectlanguage{english}
At Rome, on the Ardeatine Way, in 
 the persecution of Diocletian, the birthday of the saintly brothers Mark and 
 Marcellian, martyrs, who were arrested by the judge Fabian, tied to a stake, 
 and had sharp nails driven into their feet. Because they would not 
 cease praising the name of Christ they were pierced through the sides with 
 lances, and thus went to the kingdom of heaven with the glory of martyrdom.
\switchcolumn*
\selectlanguage{latin}
Málacæ, in Hispánia, 
 sanctórum Mártyrum Cyríaci, et Paulæ Vírginis, qui, lapídibus óbruti, inter 
 saxa cælo ánimas reddidérunt.
\switchcolumn
\selectlanguage{english}
At Malaga in Spain, the holy martyrs 
 Cyriacus and the virgin Paula, who were overwhelmed with stones, and yielded 
 up their souls to God.
\switchcolumn*
\selectlanguage{latin}
Trípoli, in Phœnícia, 
 sancti Leóntii mílitis, qui, sub Hadriáno Præside, una cum Hypátio Tribúno 
 et Theodúlo, quos convértit ad Christum, per acérba torménta pervénit ad 
 corónam martyrii.
\switchcolumn
\selectlanguage{english}
At Tripoli in Phoenicia, in the time 
 of the governor Adrian, St. Leontius, a soldier, who attained the crown of 
 martyrdom through bitter torments together with the tribune Hypatius and 
 Theodulus, whom he had converted to Christ.
\switchcolumn*
\selectlanguage{latin}
Eódem die sancti 
 Æthérii Mártyris, qui, in persecutióne Diocletiáni, post ignes et álios 
 cruciátus, gládio cæsus est.
\switchcolumn
\selectlanguage{english}
The same day, St. Aetherius, martyr, 
 in the persecution of Diocletian. After enduring fire and other 
 torments, he was put to death with the sword.
\switchcolumn*
\selectlanguage{latin}
Alexandríæ pássio 
 sanctæ Marínæ Vírginis.
\switchcolumn
\selectlanguage{english}
At Alexandria, the passion of St. 
 Marina, virgin.
\switchcolumn*
\selectlanguage{latin}
Burdígalæ sancti Amándi, 
 Epíscopi et Confessóris.
\switchcolumn
\selectlanguage{english}
At Bordeaux, St. Amandus, bishop and 
 confessor.
\switchcolumn*
\selectlanguage{latin}
Apud Saccam, in Sicília, 
 sancti Calógeri Eremítæ, cujus sánctitas in energúmenis liberándis máxime 
 effúlget.
\switchcolumn
\selectlanguage{english}
At Sacca in Sicily, St. Calogerus, 
 hermit, whose holiness is shewn especially in the deliverance of possessed 
 persons.
\switchcolumn*
\selectlanguage{latin}
Schonáugiæ, in Germánia, sanctæ Elísabeth Vírginis, ob monásticæ vitæ observántiam célebris.
\switchcolumn
\selectlanguage{english}
At Schongau in Germany, St. 
 Elizabeth, virgin, celebrated for her observance of the monastic life.
\switchcolumn*
\selectlanguage{latin}
\end{paracol}


% ---- martyrology/mart06/mart0619.htm
\needspace{10\baselineskip}
\begin{paracol}{2}
\selectlanguage{latin}
\begin{center}{\color{gregoriocolor} Tertiodécimo Kaléndas Júlii. 
 Luna\dots\ }\end{center}
\switchcolumn
\selectlanguage{english}
\begin{center}{\color{gregoriocolor} The 19\textsuperscript{th} Day of 
 June. The\dots\ Day of the Moon.}\end{center}
\end{paracol}

\noindent\begin{tabularx}{\linewidth}{*{19}{>{\centering\arraybackslash}X}}
 \textcolor{gregoriocolor}{a} & \textcolor{gregoriocolor}{b} & \textcolor{gregoriocolor}{c} & \textcolor{gregoriocolor}{d} & \textcolor{gregoriocolor}{e} & \textcolor{gregoriocolor}{f} & \textcolor{gregoriocolor}{g} & \textcolor{gregoriocolor}{h} & \textcolor{gregoriocolor}{i} & \textcolor{gregoriocolor}{k} & \textcolor{gregoriocolor}{l} & \textcolor{gregoriocolor}{m} & \textcolor{gregoriocolor}{n} & \textcolor{gregoriocolor}{p} & \textcolor{gregoriocolor}{q} & \textcolor{gregoriocolor}{r} & \textcolor{gregoriocolor}{s} & \textcolor{gregoriocolor}{t} & \textcolor{gregoriocolor}{u} \\
 23 & 24 & 25 & 26 & 27 & 28 & 29 & 1 & 2 & 3 & 4 & 5 & 6 & 7 & 8 & 9 & 10 & 11 & 12 \\
\end{tabularx}
\vspace{0.5\baselineskip}
\noindent\begin{tabularx}{\linewidth}{*{12}{>{\centering\arraybackslash}X}}
 \textcolor{gregoriocolor}{A} & \textcolor{gregoriocolor}{B} & \textcolor{gregoriocolor}{C} & \textcolor{gregoriocolor}{D} & \textcolor{gregoriocolor}{E} & F & \textcolor{gregoriocolor}{F} & \textcolor{gregoriocolor}{G} & \textcolor{gregoriocolor}{H} & \textcolor{gregoriocolor}{M} & \textcolor{gregoriocolor}{N} & \textcolor{gregoriocolor}{P} \\
 13 & 14 & 15 & 16 & 17 & 18 & 17 & 18 & 19 & 20 & 21 & 22 \\
\end{tabularx}

\begin{paracol}{2}
\selectlanguage{latin}
\lettrine[lines=2]{F}{loréntiæ} sanctæ 
 Juliánæ Falconériæ Vírginis, quæ Sorórum Ordinis Servórum beátæ Maríæ 
 Vírginis fuit Institútrix, et a Cleménte Duodécimo, Pontífice Máximo, in 
 sanctárum Vírginum númerum reláta est.
\switchcolumn
\selectlanguage{english}
\lettrine[lines=2]{A}{t} Florence, St. Juliana Falconieri, 
 virgin, foundress of the Sisters of the Order of the Servants of the Blessed 
 Virgin Mary, who was placed among the holy virgins by the Sovereign Pontiff, 
 Clement XII.
\switchcolumn*
\selectlanguage{latin}
Medioláni sanctórum 
 Mártyrum Gervásii et Protásii fratrum, ex quibus priórem támdiu jussit 
 Astásius Judex plumbátis cædi, quoúsque ille spíritum exhaláret; posteriórem 
 vero, fústibus cæsum, cápite truncári. Horum córpora, Dómino revelánte, 
 beátus Ambrósius sánguine conspérsa et ita incorrúpta réperit, ac si eo die 
 ipsi fuíssent interémpti; in quorum translatióne cæcus, féretri tactu, lumen 
 recépit, et plúrimi, vexáti a dæmónibus, liberáti sunt.
\switchcolumn
\selectlanguage{english}
At Milan, the holy martyrs Gervase 
 and Protase, brothers. The former, by order of the judge Astasius, was 
 scourged with leaded whips for so long that he expired. The latter, 
 after being scourged with rods, was beheaded. Through divine 
 revelation their bodies were found by St. Ambrose. They were partly 
 covered with blood, and as free from corruption as if they had been put to 
 death that very day. When the translation took place, a blind man 
 recovered his sight by touching their relics, and many persons possessed by 
 demons were delivered.
\switchcolumn*
\selectlanguage{latin}
In monastério Vallis 
 Castri, in Picéno, natális sancti Romuáldi Ravennátis, Anachorétæ et 
 Monachórum Camaldulénsium Patris; qui collápsam in Itália eremíticam 
 disciplínam restítuit ac mirífice propagávit. Ejus tamen festívitas 
 recólitur séptimo Idus Februárii, quo die sacræ ipsíus relíquiæ Fabriánum 
 sunt translátæ.
\switchcolumn
\selectlanguage{english}
At the monastery in the valley of 
 Castro in Piceno, the birthday of St. Romuald, anchoret, a native of 
 Ravenna. He was the founder of the Camaldolese monks, and he restored 
 and greatly extended monastic discipline, which was much relaxed in Italy. 
 His feast is observed on the 7th of February, on which day his sacred relics 
 were transferred to Fabriano.
\switchcolumn*
\selectlanguage{latin}
Arétii, in Túscia, 
 sanctórum Mártyrum Gaudéntii Epíscopi, et Culmátii Diáconi, qui, témpore 
 Valentiniáni, furóre Gentílium cæsi sunt.
\switchcolumn
\selectlanguage{english}
At Arezzo in Tuscany, the holy 
 martyrs Gaudentius, bishop, and Culmatius, deacon, who were murdered by the 
 furious heathen, during the reign of Valentinian.
\switchcolumn*
\selectlanguage{latin}
Eódem die sancti 
 Bonifátii, Epíscopi et Mártyris; qui fuit beáti Romuáldi discípulus. 
 Hic, a Gregório Quinto, Románo Pontífice, ad prædicándum Evangélium in 
 Rússiam missus, ibi, cum per ignem transísset illæsus, Regémque ac pópulum 
 baptizásset, a furénte Regis fratre necátus est, atque sic optátam martyrii 
 corónam accépit.
\switchcolumn
\selectlanguage{english}
Also, St. Boniface, martyr, a 
 disciple of blessed Romuald, who was sent by the Roman Pontiff, Gregory V, 
 to preach the Gospel in Russia. Having passed through fire uninjured, 
 and baptized the king and his people, he was killed by the enraged brother 
 of the king, and thus gained the palm of martyrdom which he ardently 
 desired.
\switchcolumn*
\selectlanguage{latin}
Ravénnæ sancti Ursicíni 
 Mártyris, qui, sub Paulíno Júdice, post plúrima torménta, in Dómini 
 confessióne pérmanens immóbilis, cápitis abscissióne martyrium complévit.
\switchcolumn
\selectlanguage{english}
At Ravenna, St. Ursicinus, martyr, 
 who remained constant through many torments in the confession of martyrdom 
 by being beheaded.
\switchcolumn*
\selectlanguage{latin}
Sozópoli, in Pisídia, 
 sancti Zósimi Mártyris, qui, in persecutióne Trajáni, sub Domitiáno Præside, 
 post acérbos cruciátus, cápite amputáto, victor migrávit ad Dóminum.
\switchcolumn
\selectlanguage{english}
At Sozopolis, under the governor 
 Domitian, during the persecution of Trajan, St. Zosimus, martyr, who 
 suffered bitter tortures, was beheaded, and thus triumphantly went to 
 heaven.
\switchcolumn*
\selectlanguage{latin}
\end{paracol}


% ---- martyrology/mart06/mart0620.htm
\needspace{10\baselineskip}
\begin{paracol}{2}
\selectlanguage{latin}
\begin{center}{\color{gregoriocolor} Duodécimo Kaléndas Júlii. 
 Luna\dots\ }\end{center}
\switchcolumn
\selectlanguage{english}
\begin{center}{\color{gregoriocolor} The 
 20\textsuperscript{th} Day of 
 June. The\dots\ Day of the Moon.}\end{center}
\end{paracol}

\noindent\begin{tabularx}{\linewidth}{*{19}{>{\centering\arraybackslash}X}}
 \textcolor{gregoriocolor}{a} & \textcolor{gregoriocolor}{b} & \textcolor{gregoriocolor}{c} & \textcolor{gregoriocolor}{d} & \textcolor{gregoriocolor}{e} & \textcolor{gregoriocolor}{f} & \textcolor{gregoriocolor}{g} & \textcolor{gregoriocolor}{h} & \textcolor{gregoriocolor}{i} & \textcolor{gregoriocolor}{k} & \textcolor{gregoriocolor}{l} & \textcolor{gregoriocolor}{m} & \textcolor{gregoriocolor}{n} & \textcolor{gregoriocolor}{p} & \textcolor{gregoriocolor}{q} & \textcolor{gregoriocolor}{r} & \textcolor{gregoriocolor}{s} & \textcolor{gregoriocolor}{t} & \textcolor{gregoriocolor}{u} \\
 24 & 25 & 26 & 27 & 28 & 29 & 1 & 2 & 3 & 4 & 5 & 6 & 7 & 8 & 9 & 10 & 11 & 12 & 13 \\
\end{tabularx}
\vspace{0.5\baselineskip}
\noindent\begin{tabularx}{\linewidth}{*{12}{>{\centering\arraybackslash}X}}
 \textcolor{gregoriocolor}{A} & \textcolor{gregoriocolor}{B} & \textcolor{gregoriocolor}{C} & \textcolor{gregoriocolor}{D} & \textcolor{gregoriocolor}{E} & F & \textcolor{gregoriocolor}{F} & \textcolor{gregoriocolor}{G} & \textcolor{gregoriocolor}{H} & \textcolor{gregoriocolor}{M} & \textcolor{gregoriocolor}{N} & \textcolor{gregoriocolor}{P} \\
 14 & 15 & 16 & 17 & 18 & 19 & 18 & 19 & 20 & 21 & 22 & 23 \\
\end{tabularx}

\begin{paracol}{2}
\selectlanguage{latin}
\lettrine[lines=2]{I}{n} ínsula Póntia 
 natális sancti Silvérii, Papæ et Mártyris, qui, cum Anthimum, Epíscopum hæréticum, a suo prædecessóre Agapíto depósitum, restitúere noluísset, a 
 Belisário, agénte ímpia Theodóra Augústa, in exsílium pulsus est, et ibídem, 
 pro fide cathólica multis ærúmnis conféctus, defécit.
\switchcolumn
\selectlanguage{english}
\lettrine[lines=2]{O}{n} the island of Pontia, the 
 birthday of St. Silverius, pope and martyr. For refusing to reinstate 
 the heretical bishop Anthimus who had been deposed by his predecessor 
 Agapitus, he was banished to the isle of Pontia by Belisarius, prompted by 
 the wicked empress Theodora. He died there, consumed by many 
 tribulations for the Catholic faith.
\switchcolumn*
\selectlanguage{latin}
Romæ deposítio sancti 
 Nováti, qui fuit fílius beáti Pudéntis Senatóris, et frater sancti Timóthei 
 Presbyteri et sanctárum Christi Vírginum Pudentiánæ et Praxédis, qui ab 
 Apóstolis erudíti sunt in fide. Ipsórum vero domus, in Ecclésiam 
 commutáta, Pastóris Títulus appellátur.
\switchcolumn
\selectlanguage{english}
At Rome, the death of St. Novatius, 
 son of the blessed senator Pudens, and brother of the saintly priest Timothy 
 and the holy virgins of Christ, Pudentiana and Praxedes, who had been 
 instructed in the faith by the apostles. Their house was converted 
 into a church, and bore the title of the Shepherd.
\switchcolumn*
\selectlanguage{latin}
Tomis, in Ponto, 
 sanctórum Mártyrum Pauli et Cyríaci.
\switchcolumn
\selectlanguage{english}
At Tomis in Pontus, the holy martyrs 
 Paul and Cyriacus.
\switchcolumn*
\selectlanguage{latin}
Petræ, in Palæstína, 
 sancti Macárii Epíscopi, qui, ab Ariánis passus multa et in Africa relegátus, 
 Conféssor quiévit in Dómino.
\switchcolumn
\selectlanguage{english}
At Petra in Palestine, St. Macarius, 
 a bishop, who suffered many things from the Arians, and was banished to 
 Africa where he rested in the Lord.
\switchcolumn*
\selectlanguage{latin}
Híspali, in Hispánia, 
 sanctæ Florentínæ Vírginis, soróris sanctórum Leándri et Isidóri Episcopórum.
\switchcolumn
\selectlanguage{english}
At Seville in Spain, the holy virgin 
 Florentina, sister of the sainted bishops Leander and Isidore.
\switchcolumn*
\selectlanguage{latin}
\end{paracol}


% ---- martyrology/mart06/mart0621.htm
\needspace{10\baselineskip}
\begin{paracol}{2}
\selectlanguage{latin}
\begin{center}{\color{gregoriocolor} Undécimo Kaléndas Júlii. 
 Luna\dots\ }\end{center}
\switchcolumn
\selectlanguage{english}
\begin{center}{\color{gregoriocolor} The 
 21\textsuperscript{st} Day of 
 June. The\dots\ Day of the Moon.}\end{center}
\end{paracol}

\noindent\begin{tabularx}{\linewidth}{*{19}{>{\centering\arraybackslash}X}}
 \textcolor{gregoriocolor}{a} & \textcolor{gregoriocolor}{b} & \textcolor{gregoriocolor}{c} & \textcolor{gregoriocolor}{d} & \textcolor{gregoriocolor}{e} & \textcolor{gregoriocolor}{f} & \textcolor{gregoriocolor}{g} & \textcolor{gregoriocolor}{h} & \textcolor{gregoriocolor}{i} & \textcolor{gregoriocolor}{k} & \textcolor{gregoriocolor}{l} & \textcolor{gregoriocolor}{m} & \textcolor{gregoriocolor}{n} & \textcolor{gregoriocolor}{p} & \textcolor{gregoriocolor}{q} & \textcolor{gregoriocolor}{r} & \textcolor{gregoriocolor}{s} & \textcolor{gregoriocolor}{t} & \textcolor{gregoriocolor}{u} \\
 25 & 26 & 27 & 28 & 29 & 1 & 2 & 3 & 4 & 5 & 6 & 7 & 8 & 9 & 10 & 11 & 12 & 13 & 14 \\
\end{tabularx}
\vspace{0.5\baselineskip}
\noindent\begin{tabularx}{\linewidth}{*{12}{>{\centering\arraybackslash}X}}
 \textcolor{gregoriocolor}{A} & \textcolor{gregoriocolor}{B} & \textcolor{gregoriocolor}{C} & \textcolor{gregoriocolor}{D} & \textcolor{gregoriocolor}{E} & F & \textcolor{gregoriocolor}{F} & \textcolor{gregoriocolor}{G} & \textcolor{gregoriocolor}{H} & \textcolor{gregoriocolor}{M} & \textcolor{gregoriocolor}{N} & \textcolor{gregoriocolor}{P} \\
 15 & 16 & 17 & 18 & 19 & 20 & 19 & 20 & 21 & 22 & 23 & 24 \\
\end{tabularx}

\begin{paracol}{2}
\selectlanguage{latin}
\lettrine[lines=2]{R}{omæ} sancti Aloísii 
 Gonzágæ, Clérici e Societáte Jesu et Confessóris, principátus contémptu et 
 innocéntia vitæ claríssimi, quem, a Summo Pontífice Benedícto Décimo tértio 
 adscríptum Sanctórum fastis et Protectórem juvénibus præsértim studiósis datum, Pius Papa Undécimus cæléstem Christiánæ juventútis univérsæ Patrónum 
 confirmávit solémniter atque íterum declarávit.
\switchcolumn
\selectlanguage{english}
\lettrine[lines=2]{A}{t} Rome, St. Aloysius Gonzaga, 
 cleric of the Society of Jesus and confessor, most renowned for his contempt 
 of the princely dignity and the innocence of his life. Pope Benedict 
 XIII placed him on the canon of the saints as special protector of young 
 students; Pope Pius XI confirmed this and again solemnly declared him to be 
 the heavenly patron of all Christian youth.
\switchcolumn*
\selectlanguage{latin}
Item Romæ sanctæ 
 Demétriæ Vírginis, quæ, sanctórum Mártyrum Flaviáni et Dafrósæ fília ac 
 sanctæ Vírginis et Mártyris Bibiánæ soror, ipsa quoque, sub Juliáno Apóstata, 
 martyrio coronáta est.
\switchcolumn
\selectlanguage{english}
Also at Rome, St. Demetria, virgin, 
 daughter of the holy martyrs Flavian and Dafrosa, and the sister of St. 
 Bibiana, virgin and martyr. She was crowned with martyrdom under 
 Julian the Apostate.
\switchcolumn*
\selectlanguage{latin}
Eódem die sancti 
 Eusébii, Samosaténi Epíscopi; qui, témpore Constántii, Imperatóris Ariáni, 
 sub hábitu militári incógnitus Ecclésias Dei visitábat, ut in fide cathólica 
 illas confirmáret. Deínde, sub Valénte, in Thráciam relegátur; sed, 
 réddita pace Ecclésiæ, témpore Theodósii, ab exsílio revocátus, tandem, cum 
 íterum Ecclésias visitáret, ei caput, tégula per mulíerem Ariánam désuper 
 immíssa, confráctum est, sicque Martyr occúbuit.
\switchcolumn
\selectlanguage{english}
The same day, St. Eusebius, bishop 
 of Samosata. In the time of the Arian emperor Constantius, he 
 disguised himself in military dress and visited the churches of God to 
 confirm them in the faith. He was banished into Thrace by Valens, but 
 when peace was restored to the Church in the reign of Theodosius, he was 
 recalled. When he again visited the churches, an Arian woman threw a 
 tile down upon him, which fractured his skull and made him a martyr.
\switchcolumn*
\selectlanguage{latin}
Icónii, in Lycaónia, 
 sancti Teréntii, Epíscopi et Mártyris.
\switchcolumn
\selectlanguage{english}
At Iconium in Lycaonia, St. Terence, 
 bishop and martyr.
\switchcolumn*
\selectlanguage{latin}
Syracúsis, in Sicília, 
 natális sanctórum Mártyrum Rufíni et Mártiæ.
\switchcolumn
\selectlanguage{english}
At Syracuse in Sicily, the birthday 
 of the holy martyrs Rufinus and Martia.
\switchcolumn*
\selectlanguage{latin}
In Africa sanctórum 
 Mártyrum Cyríaci et Apollináris.
\switchcolumn
\selectlanguage{english}
In Africa, the holy martyrs Cyriacus 
 and Apollinaris.
\switchcolumn*
\selectlanguage{latin}
Mogúntiæ sancti Albáni 
 Mártyris, qui ob Christi fidem, post longos labóres et dura certámina, 
 factus est dignus coróna vitæ.
\switchcolumn
\selectlanguage{english}
At Mainz, St. Alban, martyr, who was 
 made worthy of the crown of life, after long labors and severe combats.
\switchcolumn*
\selectlanguage{latin}
Papíæ sancti Urciscéni, 
 Epíscopi et Confessóris.
\switchcolumn
\selectlanguage{english}
At Pavia, St. Ursiscenus, bishop and 
 confessor.
\switchcolumn*
\selectlanguage{latin}
Apud Tungrénses sancti 
 Martíni Epíscopi.
\switchcolumn
\selectlanguage{english}
At Tongres, St. Martin, bishop.
\switchcolumn*
\selectlanguage{latin}
In pago Ebroicénsi 
 sancti Leutfrídi Abbátis.
\switchcolumn
\selectlanguage{english}
In the parts of Evreux, St. Leutfrid, 
 abbot.
\switchcolumn*
\selectlanguage{latin}
\end{paracol}


% ---- martyrology/mart06/mart0622.htm
\needspace{10\baselineskip}
\begin{paracol}{2}
\selectlanguage{latin}
\begin{center}{\color{gregoriocolor} Décimo Kaléndas Júlii. 
 Luna\dots\ }\end{center}
\switchcolumn
\selectlanguage{english}
\begin{center}{\color{gregoriocolor} The 
 22\textsuperscript{nd} Day of 
 June. The\dots\ Day of the Moon.}\end{center}
\end{paracol}

\noindent\begin{tabularx}{\linewidth}{*{19}{>{\centering\arraybackslash}X}}
 \textcolor{gregoriocolor}{a} & \textcolor{gregoriocolor}{b} & \textcolor{gregoriocolor}{c} & \textcolor{gregoriocolor}{d} & \textcolor{gregoriocolor}{e} & \textcolor{gregoriocolor}{f} & \textcolor{gregoriocolor}{g} & \textcolor{gregoriocolor}{h} & \textcolor{gregoriocolor}{i} & \textcolor{gregoriocolor}{k} & \textcolor{gregoriocolor}{l} & \textcolor{gregoriocolor}{m} & \textcolor{gregoriocolor}{n} & \textcolor{gregoriocolor}{p} & \textcolor{gregoriocolor}{q} & \textcolor{gregoriocolor}{r} & \textcolor{gregoriocolor}{s} & \textcolor{gregoriocolor}{t} & \textcolor{gregoriocolor}{u} \\
 26 & 27 & 28 & 29 & 1 & 2 & 3 & 4 & 5 & 6 & 7 & 8 & 9 & 10 & 11 & 12 & 13 & 14 & 15 \\
\end{tabularx}
\vspace{0.5\baselineskip}
\noindent\begin{tabularx}{\linewidth}{*{12}{>{\centering\arraybackslash}X}}
 \textcolor{gregoriocolor}{A} & \textcolor{gregoriocolor}{B} & \textcolor{gregoriocolor}{C} & \textcolor{gregoriocolor}{D} & \textcolor{gregoriocolor}{E} & F & \textcolor{gregoriocolor}{F} & \textcolor{gregoriocolor}{G} & \textcolor{gregoriocolor}{H} & \textcolor{gregoriocolor}{M} & \textcolor{gregoriocolor}{N} & \textcolor{gregoriocolor}{P} \\
 16 & 17 & 18 & 19 & 20 & 21 & 20 & 21 & 22 & 23 & 24 & 25 \\
\end{tabularx}

\begin{paracol}{2}
\selectlanguage{latin}
\lettrine[lines=2]{A}{pud} Nolam, Campániæ 
 urbem, natális beáti Paulíni, Epíscopi et Confessóris, qui ex nobilíssimo et 
 opulentíssimo factus est pro Christo pauper et húmilis, et, quod supérerat, 
 seípsum pro rediméndo víduæ fílio, quem Wándali, Campánia devastáta, 
 captívum in Africam abdúxerant, in servitútem dedit. Cláruit autem non 
 solum eruditióne et copiósa vitæ sanctitáte, sed étiam poténtia advérsus 
 dæmones; ejúsque præcláras laudes sancti Ambrósius, Hierónymus, Augustínus 
 et Gregórius Papa scriptis suis celebrárunt. Ipsíus corpus, póstea 
 Benevéntum et inde Romam translátum, tandem, Summi Pontíficis Pii Décimi 
 jussu, Nolæ restitútum fuit.
\switchcolumn
\selectlanguage{english}
\lettrine[lines=2]{A}{t} Nola in Campania, the birthday of 
 blessed Paulinus, bishop and confessor, who, although a noble and wealthy 
 man, made himself poor and humble for Christ; and what is still more 
 admirable, became a slave to liberate a widow's son who had been carried to 
 Africa by the Vandals when they devastated Campania. He was 
 celebrated, not only for his learning and great holiness of life, but also 
 for his power over demons. His great merit has been extolled by Saints 
 Ambrose, Jerome, Augustine, and Gregory in their writings. His body 
 was translated to Benevento, and later to Rome, but was taken back to Nola 
 by the order of Pope Pius X.
\switchcolumn*
\selectlanguage{latin}
Londíni in Anglia, 
 sancti Joánnis Fisher, Epíscopi Roffénsis et Cardinális, qui pro fide 
 cathólica et Románi Pontificis primátu, jubénte Henríco Octávo Rege, 
 decollátus est.
\switchcolumn
\selectlanguage{english}
At London in England, on Tower Hill, 
 St. John Fisher, bishop of Rochester and cardinal of the Holy Roman Church. 
 For the defence of the Catholic faith and the primacy of the Roman Pontiff 
 he was beheaded by order of King Henry VIII.
\switchcolumn*
\selectlanguage{latin}
In Monte Ararath pássio 
 sanctórum Mártyrum decem míllium, crucifixórum.
\switchcolumn
\selectlanguage{english}
On Mt. Ararat, the martyrdom of ten 
 thousand holy martyrs, who were crucified.
\switchcolumn*
\selectlanguage{latin}
Verolámii, in 
 Británnia, sancti Albáni Mártyris, qui, témpore Diocletiáni, pro Clérico 
 hóspite, quem domi excéperat et a quo Christiánæ fídei præceptiónibus 
 imbútus fúerat, seípsum, commutáta veste, trádidit; et hanc ob causam, post 
 vérbera et acérba torménta, cápite plexus est. Passus est étiam cum 
 illo unus de milítibus, qui, dum eum dúceret ad supplícium, in via convérsus 
 est ad Christum, et mox, gládio decollátus, próprio sánguine méruit 
 baptizári. Hoc autem nóbile sancti Albáni ac Sócii durátum pro Deo 
 certámen sanctus Beda Venerábilis descrípsit.
\switchcolumn
\selectlanguage{english}
At Verulam in England, in the time 
 of Diocletian, St. Alban, martyr, who gave himself up in order to save a 
 cleric whom he had harboured. After being scourged and subjected to 
 bitter torments, he was sentenced to capital punishment. With him also 
 suffered one of the soldiers who led him to execution, for he was converted 
 to Christ on the way and merited to be baptized in his own blood. St. 
 Venerable Bede has left an account of the noble combat of St. Alban and his 
 companion.
\switchcolumn*
\selectlanguage{latin}
Samaríæ, in Palæstína, 
 sanctórum mille quadringentórum octogínta Mártyrum, qui, sub Rege Persárum 
 Chósroa, pro Christi fide interfécti sunt.
\switchcolumn
\selectlanguage{english}
At Samaria in Palestine, fourteen 
 hundred and eighty holy martyrs, under Chosroes, king of Persia.
\switchcolumn*
\selectlanguage{latin}
Eódem die sancti Nicétæ, 
 Romantiánæ civitátis Epíscopi, doctrína sanctísque móribus clari.
\switchcolumn
\selectlanguage{english}
The same day, St. Nicaeas, bishop of 
 the town of Romatia, celebrated for his learning and holy life.
\switchcolumn*
\selectlanguage{latin}
Neápoli, in Campánia, 
 sancti Joánnis Epíscopi, quem beátus Paulínus, Epíscopus Nolánus, ad 
 cæléstia regna vocávit.
\switchcolumn
\selectlanguage{english}
At Naples in Campania, St. John, 
 bishop, who was called to the kingdom of heaven by blessed Paulinus, bishop 
 of Nola.
\switchcolumn*
\selectlanguage{latin}
In monastério 
 Cluniacénsi, in Gállia, deposítio sanctæ Consórtiæ Vírginis.
\switchcolumn
\selectlanguage{english}
In the monastery of Cluny, St. 
 Consortia, virgin.
\switchcolumn*
\selectlanguage{latin}
Romæ beáti Innocéntii 
 Papæ Quinti, ex Ordine Prædicatórum, Confessóris, qui ad tuéndam Ecclésiæ 
 libertátem et Christianórum concórdiam suávi prudéntia adlaborávit. 
 Cultum autem, ei exhíbitum, Leo Décimus tértius, Póntifex Máximus, ratum 
 hábuit et confirmávit.
\switchcolumn
\selectlanguage{english}
At Rome, blessed Pope Innocent V, 
 who laboured with mildness and prudence to maintain liberty for the Church 
 and harmony among the Christians. The veneration paid to him was 
 approved and confirmed by Pope Leo XIII.
\switchcolumn*
\selectlanguage{latin}
Item Romæ Translátio 
 sancti Flávii Cleméntis, viri Consuláris et Mártyris; qui, sanctæ Plautíllæ 
 frater ac beátæ Vírginis et Mártyris Fláviæ Domitíllæ avúnculus, a Domitiáno 
 Imperatóre, quocum Consulátum gésserat, ob Christi fidem interémptus est. 
 Ipsíus porro corpus, in Basílica sancti Cleméntis Papæ invéntum, ibídem 
 solémni pompa recónditum est.
\switchcolumn
\selectlanguage{english}
Likewise at Rome, the translation of 
 St. Flavius Clemens, exconsul and martyr, brother of St. Plautilla and uncle 
 of St. Flavia Domitilla, virgin and martyr. He was put to death for 
 the faith of Christ by Emperor Domitian. His body was found in the 
 Basilica of Pope St. Clement, and buried there with great pomp.
\switchcolumn*
\selectlanguage{latin}
\end{paracol}


% ---- martyrology/mart06/mart0623.htm
\needspace{10\baselineskip}
\begin{paracol}{2}
\selectlanguage{latin}
\begin{center}{\color{gregoriocolor} Nono Kaléndas Júlii. 
 Luna\dots\ }\end{center}
\switchcolumn
\selectlanguage{english}
\begin{center}{\color{gregoriocolor} The 
 23\textsuperscript{rd} Day of 
 June. The\dots\ Day of the Moon.}\end{center}
\end{paracol}

\noindent\begin{tabularx}{\linewidth}{*{19}{>{\centering\arraybackslash}X}}
 \textcolor{gregoriocolor}{a} & \textcolor{gregoriocolor}{b} & \textcolor{gregoriocolor}{c} & \textcolor{gregoriocolor}{d} & \textcolor{gregoriocolor}{e} & \textcolor{gregoriocolor}{f} & \textcolor{gregoriocolor}{g} & \textcolor{gregoriocolor}{h} & \textcolor{gregoriocolor}{i} & \textcolor{gregoriocolor}{k} & \textcolor{gregoriocolor}{l} & \textcolor{gregoriocolor}{m} & \textcolor{gregoriocolor}{n} & \textcolor{gregoriocolor}{p} & \textcolor{gregoriocolor}{q} & \textcolor{gregoriocolor}{r} & \textcolor{gregoriocolor}{s} & \textcolor{gregoriocolor}{t} & \textcolor{gregoriocolor}{u} \\
 27 & 28 & 29 & 1 & 2 & 3 & 4 & 5 & 6 & 7 & 8 & 9 & 10 & 11 & 12 & 13 & 14 & 15 & 16 \\
\end{tabularx}
\vspace{0.5\baselineskip}
\noindent\begin{tabularx}{\linewidth}{*{12}{>{\centering\arraybackslash}X}}
 \textcolor{gregoriocolor}{A} & \textcolor{gregoriocolor}{B} & \textcolor{gregoriocolor}{C} & \textcolor{gregoriocolor}{D} & \textcolor{gregoriocolor}{E} & F & \textcolor{gregoriocolor}{F} & \textcolor{gregoriocolor}{G} & \textcolor{gregoriocolor}{H} & \textcolor{gregoriocolor}{M} & \textcolor{gregoriocolor}{N} & \textcolor{gregoriocolor}{P} \\
 17 & 18 & 19 & 20 & 21 & 22 & 21 & 22 & 23 & 24 & 25 & 26 \\
\end{tabularx}

\begin{paracol}{2}
\selectlanguage{latin}
\lettrine[lines=1]{V}{igília} Nativitátis 
 sancti Joánnis Baptístæ.
\switchcolumn
\selectlanguage{english}
\lettrine[lines=1]{T}{he} Vigil of St. John Baptist.
\switchcolumn*
\selectlanguage{latin}
Romæ sancti Joánnis 
 Presbyteri, qui, sub Juliáno Apóstata, via Salária véteri, ante simulácrum 
 Solis decollátus est, et corpus ejus a beáto Concórdio Presbytero juxta 
 Mártyrum Concília sepúltum.
\switchcolumn
\selectlanguage{english}
At Rome, in the reign of Julian the 
 Apostate, St. John, a priest who was beheaded on the old Salarian Way before 
 an idol of the sun. His body was buried near those of other martyrs by 
 the blessed priest Concordius.
\switchcolumn*
\selectlanguage{latin}
Item Romæ sanctæ 
 Agrippínæ, Vírginis et Mártyris, quæ sub Valeriáno Imperatóre martyrium 
 consummávit. Ipsíus autem corpus, in Sicíliam translátum ac Menis 
 cónditum, multis miráculis corúscat.
\switchcolumn
\selectlanguage{english}
Also at Rome, St. Agrippina, virgin 
 and martyr, under the emperor Valerian. Her body was taken to Sicily, 
 where it works many miracles.
\switchcolumn*
\selectlanguage{latin}
Sútrii, in Túscia, 
 sancti Felícis Presbyteri, cujus os támdiu jussit Túrcius Præféctus lápide 
 contúndi, donec ipse Felix emítteret spíritum.
\switchcolumn
\selectlanguage{english}
At Sutri in Tuscany, St. Felix, 
 priest. By the command of the prefect Turcius, he was struck on the 
 mouth with a stone until he breathed no more.
\switchcolumn*
\selectlanguage{latin}
Nicomedíæ commemorátio 
 plurimórum sanctórum Mártyrum, qui, témpore Diocletiáni, in móntibus et 
 spelúncis laténtes, pro Christi nómine martyrium læto ánimo subiérunt.
\switchcolumn
\selectlanguage{english}
At Nicomedia, in the time of 
 Diocletian, the commemoration of many holy martyrs who concealed themselves 
 in mountains and caverns, but joyfully underwent martyrdom for the name of 
 Christ.
\switchcolumn*
\selectlanguage{latin}
Philadelphíæ, in 
 Arábia, sanctórum Mártyrum Zenónis, ejúsque servi Zenæ. Hic dómini sui 
 vincti caténas exósculans, eúmque rogans ut se in torméntis partícipem 
 dignarétur habére, a milítibus tentus est, et cum ipso dómino parem martyrii 
 corónam accépit.
\switchcolumn
\selectlanguage{english}
At Philadelphia in Arabia, the holy 
 martyrs Zeno and his slave Zenas. When the latter kissed the chains of 
 his master, begging to be a partner in his torments, he was arrested by the 
 soldiers, and received the crown of martyrdom with him.
\switchcolumn*
\selectlanguage{latin}
Augústæ Taurinórum 
 sancti Joséphi Cafásso, Sacerdótis, qui levítis pietáte et sciéntiæ augéndis 
 atque damnátis cápite Deo conciliándis fuit illústris, et a Pio Papa 
 Duodécimo inter sanctos Cælites adscríptus est.
\switchcolumn
\selectlanguage{english}
At Turin, St. Joseph Cafasso, 
 priest, renowned for his piety and learning, and for his work with 
 prisoners, reconciling to God those who were preparing for execution. 
 He was added to the number of the Saints by Pope Pius XII.
\switchcolumn*
\selectlanguage{latin}
In monastério Elyénsi, 
 in Británnia, sanctæ Ediltrúdis, Regínæ et Vírginis, quæ sanctitáte et 
 miráculis clara migrávit ad Dóminum. Ipsíus autem corpus, úndecim post 
 annis, invéntum est incorrúptum.
\switchcolumn
\selectlanguage{english}
In England, in the monastery of Ely, 
 St. Etheldreda, queen and virgin, who departed for heaven with a great 
 renown for sanctity and miracles. Her body was found without 
 corruption eleven years afterwards.
\switchcolumn*
\selectlanguage{latin}
\end{paracol}


% ---- martyrology/mart06/mart0624.htm
\needspace{10\baselineskip}
\begin{paracol}{2}
\selectlanguage{latin}
\begin{center}{\color{gregoriocolor} Octávo Kaléndas Júlii. 
 Luna\dots\ }\end{center}
\switchcolumn
\selectlanguage{english}
\begin{center}{\color{gregoriocolor} The 
 24\textsuperscript{th} Day of 
 June. The\dots\ Day of the Moon.}\end{center}
\end{paracol}

\noindent\begin{tabularx}{\linewidth}{*{19}{>{\centering\arraybackslash}X}}
 \textcolor{gregoriocolor}{a} & \textcolor{gregoriocolor}{b} & \textcolor{gregoriocolor}{c} & \textcolor{gregoriocolor}{d} & \textcolor{gregoriocolor}{e} & \textcolor{gregoriocolor}{f} & \textcolor{gregoriocolor}{g} & \textcolor{gregoriocolor}{h} & \textcolor{gregoriocolor}{i} & \textcolor{gregoriocolor}{k} & \textcolor{gregoriocolor}{l} & \textcolor{gregoriocolor}{m} & \textcolor{gregoriocolor}{n} & \textcolor{gregoriocolor}{p} & \textcolor{gregoriocolor}{q} & \textcolor{gregoriocolor}{r} & \textcolor{gregoriocolor}{s} & \textcolor{gregoriocolor}{t} & \textcolor{gregoriocolor}{u} \\
 28 & 29 & 1 & 2 & 3 & 4 & 5 & 6 & 7 & 8 & 9 & 10 & 11 & 12 & 13 & 14 & 15 & 16 & 17 \\
\end{tabularx}
\vspace{0.5\baselineskip}
\noindent\begin{tabularx}{\linewidth}{*{12}{>{\centering\arraybackslash}X}}
 \textcolor{gregoriocolor}{A} & \textcolor{gregoriocolor}{B} & \textcolor{gregoriocolor}{C} & \textcolor{gregoriocolor}{D} & \textcolor{gregoriocolor}{E} & F & \textcolor{gregoriocolor}{F} & \textcolor{gregoriocolor}{G} & \textcolor{gregoriocolor}{H} & \textcolor{gregoriocolor}{M} & \textcolor{gregoriocolor}{N} & \textcolor{gregoriocolor}{P} \\
 18 & 19 & 20 & 21 & 22 & 23 & 22 & 23 & 24 & 25 & 26 & 27 \\
\end{tabularx}

\begin{paracol}{2}
\selectlanguage{latin}
\lettrine[lines=2]{N}{atívitas} sancti 
 Joánnis Baptístæ, Præcursóris Dómini, ac sanctórum Zacharíæ et Elísabeth fílii, qui Spíritu Sancto replétus est adhuc in útero matris suæ.
\switchcolumn
\selectlanguage{english}
\lettrine[lines=2]{T}{he} Nativity of St. John Baptist, 
 precursor of our Lord, son of Zachary and Elizabeth, who, while yet in the 
 womb of his mother, was filled with the Holy Ghost.
\switchcolumn*
\selectlanguage{latin}
Romæ commemorátio 
 sanctórum plurimórum Mártyrum, qui a Neróne Imperatóre, ut a se incénsæ 
 Urbis ódium avérteret, calumnióse accusáti, divérso mortis génere jussi sunt 
 sævíssime intérfici. Horum síquidem álii, ferárum tergis contécti, 
 laniátibus canum expósiti sunt; álii crúcibus affíxi; aliíque incéndio 
 tráditi, ut, ubi defecísset dies, in usum noctúrni lúminis deservírent. 
 Erant hi omnes Apostolórum discípuli, et primítiæ Mártyrum, quas Romána 
 Ecclésia, fértilis ager Mártyrum, ante Apostolórum necem transmísit ad 
 Dóminum.
\switchcolumn
\selectlanguage{english}
At Rome, in the time of Nero, the 
 commemoration of many holy martyrs. Being falsely accused of having 
 set fire to the city, they were cruelly put to death in various manners by 
 the emperor's order. Some were covered with the skins of wild beasts 
 and torn to pieces by dogs, other were fastened to crosses, others again 
 were delivered to the flames to serve as torches in the night. All 
 these were disciples of the apostles, and the first fruits of the martyrs 
 which the Roman Church, a field so fertile in martyrs, offered to God even 
 before the death of the Apostles.
\switchcolumn*
\selectlanguage{latin}
Item Romæ sanctórum 
 Mártyrum Fausti et aliórum vigínti trium.
\switchcolumn
\selectlanguage{english}
In the same city, the holy martyrs 
 Faustus and twenty-three others.
\switchcolumn*
\selectlanguage{latin}
Mechlíniæ, in Brabántia, 
 pássio sancti Rumóldi, Epíscopi Dublinénsis et Mártyris, e Scotórum Rege 
 progéniti.
\switchcolumn
\selectlanguage{english}
At Mechlin in Brabant, the passion 
 of St. Rumold, bishop of Dublin and martyr. He had been the son of the 
 king of the Scots.
\switchcolumn*
\selectlanguage{latin}
Sátalis, in Arménia, sanctórum Mártyrum septem fratrum, scílicet Oréntii, Heróis, Pharnácii, 
 Firmíni, Firmi, Cyríaci et Longíni mílitum; qui a Maximiáno Imperatóre, eo 
 quod Christiáni essent, cíngulo militári priváti sunt, et, ab ínvicem 
 separáti atque in divérsa loca abdúcti, in dolóribus et ærúmnis pósiti, 
 quievérunt in Dómino.
\switchcolumn
\selectlanguage{english}
At Satalis in Armenia, seven saintly 
 brothers, all martyrs: Orentius, Heros, Pharnacius, Firminus, Firmus, 
 Cyriacus and Longinus, who owe their martyrdom to Emperor Maximian. 
 Because they were Christians, they were deprived of the military belt by his 
 command, then separated from one another, hurried away to different places, 
 and in the midst of painful trials found their repose in the Lord.
\switchcolumn*
\selectlanguage{latin}
In vico Christólio, in 
 território Parisiénsi, pássio sanctórum Mártyrum Agoárdi et Aglibérti, cum 
 áliis innúmeris promíscui sexus.
\switchcolumn
\selectlanguage{english}
In the diocese of Paris, at Creteil, 
 the martyrdom of the Saints Agoard and Aglibert, with a great multitude of 
 others of both sexes.
\switchcolumn*
\selectlanguage{latin}
Augustodúni deposítio 
 sancti Simplícii, Epíscopi et Confessóris.
\switchcolumn
\selectlanguage{english}
At Autun, the death of St. 
 Simplicius, bishop and confessor.
\switchcolumn*
\selectlanguage{latin}
Láubiis, in Bélgio, 
 sancti Theodúlphi Epíscopi.
\switchcolumn
\selectlanguage{english}
At Lobbes in Belgium, St. 
 Theodulphus, bishop.
\switchcolumn*
\selectlanguage{latin}
\end{paracol}


% ---- martyrology/mart06/mart0625.htm
\needspace{10\baselineskip}
\begin{paracol}{2}
\selectlanguage{latin}
\begin{center}{\color{gregoriocolor} Séptimo Kaléndas Júlii. 
 Luna\dots\ }\end{center}
\switchcolumn
\selectlanguage{english}
\begin{center}{\color{gregoriocolor} The 
 25\textsuperscript{th} Day of 
 June. The\dots\ Day of the Moon.}\end{center}
\end{paracol}

\noindent\begin{tabularx}{\linewidth}{*{19}{>{\centering\arraybackslash}X}}
 \textcolor{gregoriocolor}{a} & \textcolor{gregoriocolor}{b} & \textcolor{gregoriocolor}{c} & \textcolor{gregoriocolor}{d} & \textcolor{gregoriocolor}{e} & \textcolor{gregoriocolor}{f} & \textcolor{gregoriocolor}{g} & \textcolor{gregoriocolor}{h} & \textcolor{gregoriocolor}{i} & \textcolor{gregoriocolor}{k} & \textcolor{gregoriocolor}{l} & \textcolor{gregoriocolor}{m} & \textcolor{gregoriocolor}{n} & \textcolor{gregoriocolor}{p} & \textcolor{gregoriocolor}{q} & \textcolor{gregoriocolor}{r} & \textcolor{gregoriocolor}{s} & \textcolor{gregoriocolor}{t} & \textcolor{gregoriocolor}{u} \\
 29 & 1 & 2 & 3 & 4 & 5 & 6 & 7 & 8 & 9 & 10 & 11 & 12 & 13 & 14 & 15 & 16 & 17 & 18 \\
\end{tabularx}
\vspace{0.5\baselineskip}
\noindent\begin{tabularx}{\linewidth}{*{12}{>{\centering\arraybackslash}X}}
 \textcolor{gregoriocolor}{A} & \textcolor{gregoriocolor}{B} & \textcolor{gregoriocolor}{C} & \textcolor{gregoriocolor}{D} & \textcolor{gregoriocolor}{E} & F & \textcolor{gregoriocolor}{F} & \textcolor{gregoriocolor}{G} & \textcolor{gregoriocolor}{H} & \textcolor{gregoriocolor}{M} & \textcolor{gregoriocolor}{N} & \textcolor{gregoriocolor}{P} \\
 19 & 20 & 21 & 22 & 23 & 24 & 23 & 24 & 25 & 26 & 27 & 28 \\
\end{tabularx}

\begin{paracol}{2}
\selectlanguage{latin}
\lettrine[lines=2]{I}{n} território Guléti, 
 prope Nuscum, sancti Guliélmi Confessóris, Patris Eremitárum Montis Vírginis.
\switchcolumn
\selectlanguage{english}
\lettrine[lines=2]{I}{n} the territory of Guletto near 
 Nusco, St. William, confessor, founder of the hermits of Monte Vergine.
\switchcolumn*
\selectlanguage{latin}
Apud Berœam natális 
 sancti Sosípatris, qui fuit discípulus beáti Pauli Apóstoli.
\switchcolumn
\selectlanguage{english}
At Beraea, the birthday of St. 
 Sosipater, disciple of the blessed apostle Paul.
\switchcolumn*
\selectlanguage{latin}
Romæ sanctæ Lúciæ, 
 Vírginis et Mártyris, cum áliis vigínti duóbus.
\switchcolumn
\selectlanguage{english}
At Rome, St. Lucy, virgin and 
 martyr, with twenty-two others.
\switchcolumn*
\selectlanguage{latin}
Alexandríæ sancti 
 Gallicáni Mártyris, viri Consuláris, qui, sublimátus ínfulis triumphálibus 
 et Constantíno Augústo carus, a sanctis Joánne et Paulo ad Christi fidem 
 convérsus est; eáque suscépta, cum sancto Hilaríno secéssit ad Ostia 
 Tiberína, atque ibi hospitalitáti et infirmórum servítio totum se dedit. 
 Cujus rei fama in toto Orbe divulgáta, multi úndique illuc veniéntes 
 vidébant virum ex Patrício et Cónsule laváre páuperum pedes, pónere mensam, 
 aquam mánibus effúndere, languéntibus sollícite ministráre et cétera pietátis offícia exhibére. Ipse póstmodum, sub Juliána Apóstata, inde 
 expúlsus, Alexandríam perréxit; ibíque, cum a Rauciáno Júdice sacrificáre 
 cogerétur et contémneret, ideo, percússus gládio, Christi Martyr efféctus 
 est.
\switchcolumn
\selectlanguage{english}
At Alexandria, St. Gallicanus, 
 exconsul and martyr who had been honoured with a triumph, and was held in 
 affection by the emperor Constantine. Converted by Saints John and 
 Paul, he withdrew to Ostia with St. Hilarinus, and consecrated himself 
 entirely to the duties of hospitality and to the service of the sick. 
 The report of such an event spread throughout the whole world, and from all 
 sides many people came to see a man who had been a senator and consul now 
 washing the feet of the poor, preparing their table, serving them, carefully 
 waiting on the infirm, and exercising other works of mercy. Driven 
 from this place by Julian the Apostate, he repaired to Alexandria, where, 
 for refusing to sacrifice to idols, at the command of the judge Raucian, he 
 was put to the sword, and thus became a martyr of Christ.
\switchcolumn*
\selectlanguage{latin}
Sibápoli, in 
 Mesopotámia, sanctæ Febróniæ, Vírginis et Mártyris; quæ, in persecutióne 
 Diocletiáni, sub Siléno Júdice, ob fidem et pudicítiam servándam, primo 
 virgis cæsa et equúleo torta, deínde pectínibus laniáta atque igne succénsa, 
 demum, excússis déntibus ac mammis pedibúsque abscíssis, cápitis damnáta, 
 tot passiónum ornáta monílibus migrávit ad Sponsum.
\switchcolumn
\selectlanguage{english}
At Sibapolis in Syria, under the 
 governor Silenus, in the persecution of Diocletian, St. Febronia, virgin and 
 martyr. She was scourged and racked for defending her faith and her 
 chastity, then torn with iron combs and exposed to fire. Finally her 
 teeth were broken out, her breasts and feet cut away, and she was condemned 
 to capital punishment, going to her Spouse adorned with sufferings as with 
 so many jewels.
\switchcolumn*
\selectlanguage{latin}
Apud Rhégium sancti 
 Prósperi Aquitáni, ejúsdem urbis Epíscopi, eruditióne ac pietáte insígnis, 
 qui advérsus Pelagiános, pro fide cathólica, strénue decertávit.
\switchcolumn
\selectlanguage{english}
At Reggio, St. Prosper of Aquitaine, 
 bishop of that city, distinguished by his learning and piety. He 
 valiantly combated the Pelagians in defence of the Catholic faith.
\switchcolumn*
\selectlanguage{latin}
Tauríni natális sancti 
 Máximi, Epíscopi et Confessóris, doctrína et sanctitáte celebérrimi.
\switchcolumn
\selectlanguage{english}
At Turin, the birthday of St. 
 Maximus, bishop and confessor, most celebrated for his sanctity and 
 scholarship.
\switchcolumn*
\selectlanguage{latin}
In Hollándia sancti 
 Adelbérti Confessóris, qui fuit discípulus sancti Willibrórdi Epíscopi.
\switchcolumn
\selectlanguage{english}
In Holland, St. Adalbert, confessor, 
 disciple of St. Willibrord, bishop.
\switchcolumn*
\selectlanguage{latin}
\end{paracol}


% ---- martyrology/mart06/mart0626.htm
\needspace{10\baselineskip}
\begin{paracol}{2}
\selectlanguage{latin}
\begin{center}{\color{gregoriocolor} Sexto Kaléndas Júlii. 
 Luna\dots\ }\end{center}
\switchcolumn
\selectlanguage{english}
\begin{center}{\color{gregoriocolor} The 
 26\textsuperscript{th} Day of 
 June. The\dots\ Day of the Moon.}\end{center}
\end{paracol}

\noindent\begin{tabularx}{\linewidth}{*{19}{>{\centering\arraybackslash}X}}
 \textcolor{gregoriocolor}{a} & \textcolor{gregoriocolor}{b} & \textcolor{gregoriocolor}{c} & \textcolor{gregoriocolor}{d} & \textcolor{gregoriocolor}{e} & \textcolor{gregoriocolor}{f} & \textcolor{gregoriocolor}{g} & \textcolor{gregoriocolor}{h} & \textcolor{gregoriocolor}{i} & \textcolor{gregoriocolor}{k} & \textcolor{gregoriocolor}{l} & \textcolor{gregoriocolor}{m} & \textcolor{gregoriocolor}{n} & \textcolor{gregoriocolor}{p} & \textcolor{gregoriocolor}{q} & \textcolor{gregoriocolor}{r} & \textcolor{gregoriocolor}{s} & \textcolor{gregoriocolor}{t} & \textcolor{gregoriocolor}{u} \\
 1 & 2 & 3 & 4 & 5 & 6 & 7 & 8 & 9 & 10 & 11 & 12 & 13 & 14 & 15 & 16 & 17 & 18 & 19 \\
\end{tabularx}
\vspace{0.5\baselineskip}
\noindent\begin{tabularx}{\linewidth}{*{12}{>{\centering\arraybackslash}X}}
 \textcolor{gregoriocolor}{A} & \textcolor{gregoriocolor}{B} & \textcolor{gregoriocolor}{C} & \textcolor{gregoriocolor}{D} & \textcolor{gregoriocolor}{E} & F & \textcolor{gregoriocolor}{F} & \textcolor{gregoriocolor}{G} & \textcolor{gregoriocolor}{H} & \textcolor{gregoriocolor}{M} & \textcolor{gregoriocolor}{N} & \textcolor{gregoriocolor}{P} \\
 20 & 21 & 22 & 23 & 24 & 25 & 24 & 25 & 26 & 27 & 28 & 29 \\
\end{tabularx}

\begin{paracol}{2}
\selectlanguage{latin}
\lettrine[lines=2]{R}{omæ,} in monte Cælio, 
 sanctórum Mártyrum Joánnis et Pauli fratrum, quorum primus erat præpósitus 
 domus, secúndus primicérius Constántiæ Vírginis, fíliæ Constantíni 
 Imperatóris, et ambo póstea, sub Juliáno Apóstata, martyrii palmam, cædénte 
 gládio, percepérunt.
\switchcolumn
\selectlanguage{english}
\lettrine[lines=2]{A}{t} Rome on Mt. Coelius, the holy 
 martyrs John and Paul, brothers. The former was steward, the other 
 secretary of the virgin Constantia, daughter of Emperor Constantine. 
 Afterwards, under Julian the Apostate, they received the palm of martyrdom 
 by being beheaded.
\switchcolumn*
\selectlanguage{latin}
Apud Tridéntum sancti 
 Vigílii Epíscopi, qui, cum relíquias idololatríæ pénitus exstirpáre 
 conarétur, ideo, pro Christi nómine, a feris et bárbaris homínibus percússus 
 lápidum imbre, martyrium implévit.
\switchcolumn
\selectlanguage{english}
At Trent, St. Vigilius, bishop, who, 
 while he endeavoured to root out the remains of idolatry, was overwhelmed 
 with a shower of stones by cruel and barbarous men, and thus endured 
 martyrdom for the name of Christ.
\switchcolumn*
\selectlanguage{latin}
Apud Valencénas, in 
 Gállia, pássio sanctórum Mártyrum Sálvii, qui éxstitit Engolisménsis 
 Epíscopus, et Supérii.
\switchcolumn
\selectlanguage{english}
At Valenciennes, the holy martyrs Salvius, bishop of Angoulême, and Superius.
\switchcolumn*
\selectlanguage{latin}
Córdubæ, in Hispánia, 
 natális sancti Pelágii adolescéntuli, qui, ob confessiónem fídei, Regis 
 Saracenórum Abdaraméni jussu forcípibus férreis membrátim præcísus, 
 martyrium suum glorióse consummávit.
\switchcolumn
\selectlanguage{english}
At Cordova in Spain, under the 
 Saracen king Abderaliman, the birthday of St. Pelagius, a young man who 
 gloriously completed his martyrdom for the faith by having his flesh torn to 
 pieces with iron pincers.
\switchcolumn*
\selectlanguage{latin}
Bellícii, in Gállia, 
 sancti Anthélmi, qui ex majóris Carthúsiæ Prióre, factus est ejúsdem 
 civitátis Epíscopus.
\switchcolumn
\selectlanguage{english}
At Belley in France, St. Anthelmus, 
 prior of the Grande Chartreuse, who became bishop of that city.
\switchcolumn*
\selectlanguage{latin}
In pago Pictaviénsi 
 sancti Maxéntii, Presbyteri et Confessóris; qui miráculis cláruit.
\switchcolumn
\selectlanguage{english}
In the country of Poitiers, St. 
 Maxentius, priest and confessor, renowned for miracles.
\switchcolumn*
\selectlanguage{latin}
Thessalonícæ sancti 
 David Eremítæ.
\switchcolumn
\selectlanguage{english}
At Thessalonica, St. David, hermit.
\switchcolumn*
\selectlanguage{latin}
Eódem die sanctæ 
 Perseverándæ Vírginis.
\switchcolumn
\selectlanguage{english}
The same day, St. Perseveranda, 
 virgin.
\switchcolumn*
\selectlanguage{latin}
\end{paracol}


% ---- martyrology/mart06/mart0627.htm
\needspace{10\baselineskip}
\begin{paracol}{2}
\selectlanguage{latin}
\begin{center}{\color{gregoriocolor} Quinto Kaléndas Júlii. 
 Luna\dots\ }\end{center}
\switchcolumn
\selectlanguage{english}
\begin{center}{\color{gregoriocolor} The 
 27\textsuperscript{th} Day of 
 June. The\dots\ Day of the Moon.}\end{center}
\end{paracol}

\noindent\begin{tabularx}{\linewidth}{*{19}{>{\centering\arraybackslash}X}}
 \textcolor{gregoriocolor}{a} & \textcolor{gregoriocolor}{b} & \textcolor{gregoriocolor}{c} & \textcolor{gregoriocolor}{d} & \textcolor{gregoriocolor}{e} & \textcolor{gregoriocolor}{f} & \textcolor{gregoriocolor}{g} & \textcolor{gregoriocolor}{h} & \textcolor{gregoriocolor}{i} & \textcolor{gregoriocolor}{k} & \textcolor{gregoriocolor}{l} & \textcolor{gregoriocolor}{m} & \textcolor{gregoriocolor}{n} & \textcolor{gregoriocolor}{p} & \textcolor{gregoriocolor}{q} & \textcolor{gregoriocolor}{r} & \textcolor{gregoriocolor}{s} & \textcolor{gregoriocolor}{t} & \textcolor{gregoriocolor}{u} \\
 2 & 3 & 4 & 5 & 6 & 7 & 8 & 9 & 10 & 11 & 12 & 13 & 14 & 15 & 16 & 17 & 18 & 19 & 20 \\
\end{tabularx}
\vspace{0.5\baselineskip}
\noindent\begin{tabularx}{\linewidth}{*{12}{>{\centering\arraybackslash}X}}
 \textcolor{gregoriocolor}{A} & \textcolor{gregoriocolor}{B} & \textcolor{gregoriocolor}{C} & \textcolor{gregoriocolor}{D} & \textcolor{gregoriocolor}{E} & F & \textcolor{gregoriocolor}{F} & \textcolor{gregoriocolor}{G} & \textcolor{gregoriocolor}{H} & \textcolor{gregoriocolor}{M} & \textcolor{gregoriocolor}{N} & \textcolor{gregoriocolor}{P} \\
 21 & 22 & 23 & 24 & 25 & 26 & 25 & 26 & 27 & 28 & 29 & 1 \\
\end{tabularx}

\begin{paracol}{2}
\selectlanguage{latin}
\lettrine[lines=2]{I}{n} Galátia sancti 
 Crescéntis, qui fuit beáti Pauli Apóstoli discípulus. Hic, in Gállias 
 tránsitum fáciens, verbo prædicatiónis multos ad fidem Christi convértit; 
 regréssus vero ad gentem cui speciáliter datus erat Epíscopus, demum, sub 
 Trajáno, cum Gálatas ipsos usque ad finem vitæ suæ in ópere Dómini 
 confirmásset, martyrium consummávit.
\switchcolumn
\selectlanguage{english}
\lettrine[lines=2]{I}{n} Galatia, St. Crescens, disciple 
 of the blessed apostle Paul. In passing through Gaul he converted many 
 to the Christian faith by his preaching. Returning to the people for 
 whom he had been especially made bishop, he confirmed the Galatians in the 
 service of the Lord to the end of his life. He finally completed his 
 martyrdom under Trajan.
\switchcolumn*
\selectlanguage{latin}
Córdubæ, in Hispánia, 
 sanctórum Mártyrum Zóili, et aliórum decem et novem.
\switchcolumn
\selectlanguage{english}
At Cordova in Spain, St. Zoilus and 
 nineteen other martyrs.
\switchcolumn*
\selectlanguage{latin}
Cæsaréæ, in Palæstína, 
 sancti Anécti Mártyris, qui, in persecutióne Diocletiáni, sub Urbáno Præside, 
 cum álios ad martyrium hortátus esset et idóla oratióne sua prostrásset, a 
 decem milítibus verberári jussus est; ac tandem, mánibus pedibúsque præcísis, 
 truncátus cápite, martyrii corónam accépit.
\switchcolumn
\selectlanguage{english}
At Caesarea in Palestine, in the 
 persecution of Diocletian, under the governor Urban, St. Anectus, martyr. 
 For having exhorted others to suffer martyrdom, and having overthrown idols 
 by his prayers, he was scourged by ten soldiers, had his hands and feet cut 
 off, and merited the crown of martyrdom by beheading.
\switchcolumn*
\selectlanguage{latin}
Constantinópoli sancti 
 Sampsónis Presbyteri, páuperum exceptóris.
\switchcolumn
\selectlanguage{english}
At Constantinople, St. Sampson, a 
 priest, who harboured the poor.
\switchcolumn*
\selectlanguage{latin}
In castro Cainóne, in 
 Gállia, sancti Joánnis, Presbyteri et Confessóris.
\switchcolumn
\selectlanguage{english}
In the town of Chinon in France, St. 
 John, priest and confessor.
\switchcolumn*
\selectlanguage{latin}
Varadíni, in Hungária, 
 sancti Ladislái Regis, qui claríssimis miráculis usque ad diem hodiérnum 
 corúscat.
\switchcolumn
\selectlanguage{english}
At Grosswardein in Hungary, the holy 
 king Ladislaus, greatly renowned for his miracles even to this day.
\switchcolumn*
\selectlanguage{latin}
\end{paracol}


% ---- martyrology/mart06/mart0628.htm
\needspace{10\baselineskip}
\begin{paracol}{2}
\selectlanguage{latin}
\begin{center}{\color{gregoriocolor} Quarto Kaléndas Júlii. 
 Luna\dots\ }\end{center}
\switchcolumn
\selectlanguage{english}
\begin{center}{\color{gregoriocolor} The 
 28\textsuperscript{th} Day of 
 June. The\dots\ Day of the Moon.}\end{center}
\end{paracol}

\noindent\begin{tabularx}{\linewidth}{*{19}{>{\centering\arraybackslash}X}}
 \textcolor{gregoriocolor}{a} & \textcolor{gregoriocolor}{b} & \textcolor{gregoriocolor}{c} & \textcolor{gregoriocolor}{d} & \textcolor{gregoriocolor}{e} & \textcolor{gregoriocolor}{f} & \textcolor{gregoriocolor}{g} & \textcolor{gregoriocolor}{h} & \textcolor{gregoriocolor}{i} & \textcolor{gregoriocolor}{k} & \textcolor{gregoriocolor}{l} & \textcolor{gregoriocolor}{m} & \textcolor{gregoriocolor}{n} & \textcolor{gregoriocolor}{p} & \textcolor{gregoriocolor}{q} & \textcolor{gregoriocolor}{r} & \textcolor{gregoriocolor}{s} & \textcolor{gregoriocolor}{t} & \textcolor{gregoriocolor}{u} \\
 3 & 4 & 5 & 6 & 7 & 8 & 9 & 10 & 11 & 12 & 13 & 14 & 15 & 16 & 17 & 18 & 19 & 20 & 21 \\
\end{tabularx}
\vspace{0.5\baselineskip}
\noindent\begin{tabularx}{\linewidth}{*{12}{>{\centering\arraybackslash}X}}
 \textcolor{gregoriocolor}{A} & \textcolor{gregoriocolor}{B} & \textcolor{gregoriocolor}{C} & \textcolor{gregoriocolor}{D} & \textcolor{gregoriocolor}{E} & F & \textcolor{gregoriocolor}{F} & \textcolor{gregoriocolor}{G} & \textcolor{gregoriocolor}{H} & \textcolor{gregoriocolor}{M} & \textcolor{gregoriocolor}{N} & \textcolor{gregoriocolor}{P} \\
 22 & 23 & 24 & 25 & 26 & 27 & 26 & 27 & 28 & 29 & 1 & 2 \\
\end{tabularx}

\begin{paracol}{2}
\selectlanguage{latin}
\lettrine[lines=1]{V}{igília} sanctórum 
 Apostolórum Petri et Pauli.
\switchcolumn
\selectlanguage{english}
\lettrine[lines=1]{T}{he} vigil of the holy apostles Peter 
 and Paul.
\switchcolumn*
\selectlanguage{latin}
Lugdúni, in Gállia, 
 sancti Irenæi, Epíscopi et Mártyris; qui (ut scribit sanctus Hierónymus) 
 beáti Polycárpi, Smyrnénsis Epíscopi, discípulus fuit, et Apostolicórum 
 témporum vicínis. Is, cum advérsus hæréticos verbis ac scriptis 
 plúrimum decertásset, tandem, in persecutióne Sevéri, cum omni fere 
 civitátis suæ pópulo, coronátus est glorióso martyrio.
\switchcolumn
\selectlanguage{english}
At Lyons in France, St. Irenaeus, 
 bishop and martyr. St. Jerome relates that he was the disciple of 
 blessed Polycarp, bishop of Smyrna, and lived near the time of the apostles. 
 After having strenuously opposed the heretics by word and by writing, he was 
 crowned with a glorious martyrdom along with almost all the people of his 
 city, during the persecution of Severus.
\switchcolumn*
\selectlanguage{latin}
Trajécti sancti Benígni, 
 Epíscopi et Mártyris.
\switchcolumn
\selectlanguage{english}
At Utrecht, St. Benignus, bishop and 
 martyr.
\switchcolumn*
\selectlanguage{latin}
Alexandríæ, in 
 persecutióne Sevéri, sanctórum Mártyrum Plutárchi, Seréni, Heraclídis 
 catechúmeni, Herónis neóphyti, et altérius Seréni, Rháidis catechúmenæ, et 
 Potamiœnæ cum ipsíus matre Marcélla; inter quos præcípue enítuit Potamiœna 
 Virgo, quæ, primo imménsos innumerósque agónes pro virginitáte desudávit, 
 deínde étiam exquisíta et inaudíta torménta pro fide sustínuit, ad últimum, 
 simul cum matre, igne consúmpta est.
\switchcolumn
\selectlanguage{english}
At Alexandria, in the persecution of 
 Severus, the holy martyrs Plutarch, Serenus, Heraclides, catechumen, Heron, 
 a neophyte, another Serenus, Rhais, a catechumen, Potamioena and Marcella 
 her mother. Among them the virgin Potamioena is particularly 
 distinguished. She first endured many painful trials for the 
 preservation of her virginity, and then cruel and unheard-of torments for 
 the faith, after which both she and her mother were consumed with fire.
\switchcolumn*
\selectlanguage{latin}
Eódem die sancti Pápii 
 Mártyris, qui, in persecutióne Diocletiáni Imperatóris, flagris cæsus, atque 
 in lebétem, óleo et ádipe fervénti plenum, immíssus, áliaque horrénda 
 supplícia perpéssus, demum, datis cervícibus, coronátur.
\switchcolumn
\selectlanguage{english}
Also during the persecution of 
 Diocletian, St. Papius, martyr, who was scourged with knotted cords, cast 
 into a cauldron of seething oil and grease, and after other horrible 
 torments was beheaded, and thus won an eternal crown.
\switchcolumn*
\selectlanguage{latin}
Córdubæ, in Hispánia, 
 sancti Argymíri, Mónachi et Mártyris, qui in persecutióne Arábica, pro 
 Christi fide, in equúleo pósitus et ense transfóssus est.
\switchcolumn
\selectlanguage{english}
At Cordova in Spain, St. Argymirus, 
 monk and martyr, who was slain for the faith of Christ during the 
 persecution of the Arabs.
\switchcolumn*
\selectlanguage{latin}
Romæ sancti Pauli Primi, 
 Papæ et Confessóris.
\switchcolumn
\selectlanguage{english}
At Rome, Pope St. Paul I, confessor.
\switchcolumn*
\selectlanguage{latin}
Lúere, in diœcési 
 Brixiénsi, sanctæ Vincéntiæ Gerósa Vírginis, Institúti Sorórum a Caritáte 
 una cum sancta Bartholomǽa Capitánio Fundatrícis, quam Pius Papa Duodécimus 
 albo sanctárum Vírginum accénsuit.
\switchcolumn
\selectlanguage{english}
At Lovere, in the diocese of 
 Brescia, St. Vincenza Gerosa, virgin, who co-founded the Institute of the 
 Sisters of Charity with St. Bartolomea Capitanio, and whom Pope Pius XII 
 added to the list of holy virgins.
\switchcolumn*
\selectlanguage{latin}
\end{paracol}


% ---- martyrology/mart06/mart0629.htm
\needspace{10\baselineskip}
\begin{paracol}{2}
\selectlanguage{latin}
\begin{center}{\color{gregoriocolor} Tértio Kaléndas Júlii. 
 Luna\dots\ }\end{center}
\switchcolumn
\selectlanguage{english}
\begin{center}{\color{gregoriocolor} The 
 29\textsuperscript{th} Day of 
 June. The\dots\ Day of the Moon.}\end{center}
\end{paracol}

\noindent\begin{tabularx}{\linewidth}{*{19}{>{\centering\arraybackslash}X}}
 \textcolor{gregoriocolor}{a} & \textcolor{gregoriocolor}{b} & \textcolor{gregoriocolor}{c} & \textcolor{gregoriocolor}{d} & \textcolor{gregoriocolor}{e} & \textcolor{gregoriocolor}{f} & \textcolor{gregoriocolor}{g} & \textcolor{gregoriocolor}{h} & \textcolor{gregoriocolor}{i} & \textcolor{gregoriocolor}{k} & \textcolor{gregoriocolor}{l} & \textcolor{gregoriocolor}{m} & \textcolor{gregoriocolor}{n} & \textcolor{gregoriocolor}{p} & \textcolor{gregoriocolor}{q} & \textcolor{gregoriocolor}{r} & \textcolor{gregoriocolor}{s} & \textcolor{gregoriocolor}{t} & \textcolor{gregoriocolor}{u} \\
 4 & 5 & 6 & 7 & 8 & 9 & 10 & 11 & 12 & 13 & 14 & 15 & 16 & 17 & 18 & 19 & 20 & 21 & 22 \\
\end{tabularx}
\vspace{0.5\baselineskip}
\noindent\begin{tabularx}{\linewidth}{*{12}{>{\centering\arraybackslash}X}}
 \textcolor{gregoriocolor}{A} & \textcolor{gregoriocolor}{B} & \textcolor{gregoriocolor}{C} & \textcolor{gregoriocolor}{D} & \textcolor{gregoriocolor}{E} & F & \textcolor{gregoriocolor}{F} & \textcolor{gregoriocolor}{G} & \textcolor{gregoriocolor}{H} & \textcolor{gregoriocolor}{M} & \textcolor{gregoriocolor}{N} & \textcolor{gregoriocolor}{P} \\
 23 & 24 & 25 & 26 & 27 & 28 & 27 & 28 & 29 & 1 & 2 & 3 \\
\end{tabularx}

\begin{paracol}{2}
\selectlanguage{latin}
\lettrine[lines=2]{R}{omæ} natális sanctórum 
 Apostolórum Petri et Pauli, qui eódem anno eodémque die passi sunt, sub 
 Neróne Imperatóre. Horum prior, in eádem Urbe, cápite ad terram verso 
 cruci affíxus, et in Vaticáno juxta viam Triumphálem sepúltus, totíus Orbis 
 veneratióne celebrátur; postérior autem, gládio animadvérsus, et via 
 Ostiénsi sepúltus, pari honóre habétur.
\switchcolumn
\selectlanguage{english}
\lettrine[lines=2]{A}{t} Rome, the birthday of the holy 
 apostles Peter and Paul, who suffered martyrdom on the same day, under 
 Emperor Nero. Within the city the former was crucified with his head 
 downwards, and buried in the Vatican, near the Triumphal Way, where he is 
 venerated by the whole world. The latter was put to the sword and 
 buried on the Ostian Way, where he received similar honours.
\switchcolumn*
\selectlanguage{latin}
In Cypro sanctæ Maríæ, 
 matris Joánnis qui cognominátus est Marcus.
\switchcolumn
\selectlanguage{english}
In Cyprus, St. Mary, mother of John, 
 surnamed Mark.
\switchcolumn*
\selectlanguage{latin}
In castro Argentómacho, 
 in Gállia, sancti Marcélli Mártyris, qui, pro fide Christi, una cum 
 Anastásio, viro militári, cápite plexus est.
\switchcolumn
\selectlanguage{english}
At Argenton in France, St. 
 Marcellus, martyr, who was beheaded for the faith of Christ together with 
 the soldier Anastasius.
\switchcolumn*
\selectlanguage{latin}
Génuæ natális sancti 
 Syri Epíscopi.
\switchcolumn
\selectlanguage{english}
At Genoa, the birthday of St. Syrius, 
 bishop.
\switchcolumn*
\selectlanguage{latin}
Nárniæ sancti Cássii, 
 ejúsdem cívitiátis Epíscopi, de quo sanctus Gregórius Papa refert quod 
 nullus ferme dies vitæ ejus abscedébat, quo placatiónis hóstias omnipoténti 
 Deo non offérret; cui et concordábat vita, quia, cuncta quæ habébat, in 
 eleemósynis tríbuens, in hora sacrifícii totus in lácrimis diffluébat. 
 Demum, natalítio Apostolórum die, quo síngulis annis Romam veníre 
 consuéverat, in eádem Narniénsi urbe, cum Missárum solémnia celebrásset, et 
 corpus Domínicum pacémque ómnibus dedísset, migrávit ad Dóminum.
\switchcolumn
\selectlanguage{english}
At Narni, St. Cassius, bishop of 
 that city. St. Gregory relates that he permitted scarcely any day of 
 his life to pass without offering the Victim of propitiation to Almighty 
 God. It was in character with his life for he distributed in alms all 
 he possessed, and his devotion was such that abundant tears flowed from his 
 eyes during the holy Sacrifice. At last, coming to Rome on the 
 birthday of the apostles, as was his yearly custom, after having solemnly 
 celebrated Mass and given the Lord's Body and the kiss of peace to all, he 
 departed for heaven.
\switchcolumn*
\selectlanguage{latin}
In território Senonénsi 
 sanctæ Benedíctæ Vírginis.
\switchcolumn
\selectlanguage{english}
In the territory of Sens, St. 
 Benedicta, virgin.
\switchcolumn*
\selectlanguage{latin}
\end{paracol}


% ---- martyrology/mart06/mart0630.htm
\needspace{10\baselineskip}
\begin{paracol}{2}
\selectlanguage{latin}
\begin{center}{\color{gregoriocolor} Prídie Kaléndas Júlii. 
 Luna\dots\ }\end{center}
\switchcolumn
\selectlanguage{english}
\begin{center}{\color{gregoriocolor} The 
 30\textsuperscript{th} Day of 
 June. The\dots\ Day of the Moon.}\end{center}
\end{paracol}

\noindent\begin{tabularx}{\linewidth}{*{19}{>{\centering\arraybackslash}X}}
 \textcolor{gregoriocolor}{a} & \textcolor{gregoriocolor}{b} & \textcolor{gregoriocolor}{c} & \textcolor{gregoriocolor}{d} & \textcolor{gregoriocolor}{e} & \textcolor{gregoriocolor}{f} & \textcolor{gregoriocolor}{g} & \textcolor{gregoriocolor}{h} & \textcolor{gregoriocolor}{i} & \textcolor{gregoriocolor}{k} & \textcolor{gregoriocolor}{l} & \textcolor{gregoriocolor}{m} & \textcolor{gregoriocolor}{n} & \textcolor{gregoriocolor}{p} & \textcolor{gregoriocolor}{q} & \textcolor{gregoriocolor}{r} & \textcolor{gregoriocolor}{s} & \textcolor{gregoriocolor}{t} & \textcolor{gregoriocolor}{u} \\
 5 & 6 & 7 & 8 & 9 & 10 & 11 & 12 & 13 & 14 & 15 & 16 & 17 & 18 & 19 & 20 & 21 & 22 & 23 \\
\end{tabularx}
\vspace{0.5\baselineskip}
\noindent\begin{tabularx}{\linewidth}{*{12}{>{\centering\arraybackslash}X}}
 \textcolor{gregoriocolor}{A} & \textcolor{gregoriocolor}{B} & \textcolor{gregoriocolor}{C} & \textcolor{gregoriocolor}{D} & \textcolor{gregoriocolor}{E} & F & \textcolor{gregoriocolor}{F} & \textcolor{gregoriocolor}{G} & \textcolor{gregoriocolor}{H} & \textcolor{gregoriocolor}{M} & \textcolor{gregoriocolor}{N} & \textcolor{gregoriocolor}{P} \\
 24 & 25 & 26 & 27 & 28 & 29 & 28 & 29 & 1 & 2 & 3 & 4 \\
\end{tabularx}

\begin{paracol}{2}
\selectlanguage{latin}
\lettrine[lines=2]{C}{ommemorátio} sancti 
 Pauli Apóstoli.
\switchcolumn
\selectlanguage{english}
\lettrine[lines=2]{T}{he} commemoration of the holy 
 apostle Paul.
\switchcolumn*
\selectlanguage{latin}
Romæ sanctæ Lucínæ, 
 Apostolórum discípulæ, quæ, de facultátibus suis commúnicans Sanctórum 
 necessitátibus, Christiános in cárcere deténtos visitábat, ac sepultúræ 
 Mártyrum inserviébat; juxta quos et ipsa, in crypta a se constrúcta, sepúlta 
 est.
\switchcolumn
\selectlanguage{english}
At Rome, St. Lucina, a disciple of 
 the apostles, who relieved the necessities of the saints with her goods, 
 visited the Christians detained in prison, buried the martyrs, and was laid 
 by their side in a crypt which she herself had constructed.
\switchcolumn*
\selectlanguage{latin}
Item Romæ sanctæ 
 Æmiliánæ Mártyris.
\switchcolumn
\selectlanguage{english}
In the same city, St. Aemiliana, 
 martyr.
\switchcolumn*
\selectlanguage{latin}
Ipso die sanctórum 
 Mártyrum Caji Presbyteri, et Leónis Subdiáconi.
\switchcolumn
\selectlanguage{english}
The same day, the saints Caius, 
 priest, and Leo, subdeacon.
\switchcolumn*
\selectlanguage{latin}
Alexandríæ pássio 
 sancti Basílidis, qui, sub Sevéro Imperatóre, cum sanctam Potamiœnam 
 Vírginem, quam ad supplícium ducébat, ab impudicórum hóminum petulántia 
 tutátus esset, religiósi offícii mercédem ab ea recépit; nam ipsa, post 
 tríduum illi appárens et ejus cápiti corónam impónens, non tantum convértit 
 eum ad Christum, sed étiam, brevi agóne certántem, suis précibus Mártyrem 
 gloriósum effécit.
\switchcolumn
\selectlanguage{english}
At Alexandria, the passion of St. 
 Basilides, under Emperor Severus. He protected the saintly virgin 
 Potamioena from the insults of shameless men when he was leading her to 
 execution. He was rewarded for his considerate action, for at the end 
 of three days she appeared to him, placed a crown on his head, not only 
 converting him to Christ, but by her prayers making him, after a short 
 combat, a glorious martyr.
\switchcolumn*
\selectlanguage{latin}
Lemóvicis, in Aquitánia, 
 sancti Martiális Epíscopi, cum duóbus Presbyteris Alpiniáno et 
 Austricliniáno; quorum vita signis miraculórum ádmodum effúlsit.
\switchcolumn
\selectlanguage{english}
At Limoges in France, St. Martial, 
 bishop, and two priests Alpinian and Austriclinian, whose lives were 
 distinguished for miracles.
\switchcolumn*
\selectlanguage{latin}
In território 
 Vivariénsi, in Gálliis, sancti Ostiáni, Presbyteri et Confessóris.
\switchcolumn
\selectlanguage{english}
In the territory of Vivers, St. 
 Ostian, priest and confessor.
\switchcolumn*
\selectlanguage{latin}
Salánicæ, in território 
 Vicentíno, sancti Theobáldi, Presbyteri et Eremítæ, ex Campániæ Gállicæ 
 Comítibus; quem Alexánder Papa Tértius, ob sanctitátis et miraculórum famam, 
 Sanctórum cánoni adjúnxit.
\switchcolumn
\selectlanguage{english}
At Salanica, in the district of 
 Vicenza, St. Theobald, priest and hermit, one of the counts of Champagne. 
 He was added to the number of the saints by Alexander III because of his 
 holiness and miracles.
\switchcolumn*
\selectlanguage{latin}
\end{paracol}
